%% Motivic homotopy theory, as introduced by Morel and Voevodsky \cite{MorelVoevodsky}, has been a tremendously successful application of abstract homotopy theory to concrete questions of algebraic geometry and number theory. \todo{This sentence is stolen from Gheorghe--Wang--Xu. It should be rewritten.}

A striking surprise of modern computational stable homotopy theory is that the category $\SH(\C)$ of $\C$-motivic spectra, as introduced by Morel and Voevodsky \cite{MorelVoevodsky},
has found use not only in algebraic geometry, but also in the computation of stable homotopy groups of spheres, a question of purely topological origin.  In particular, the $p$-completed bi-graded homotopy groups of the unit in $\SH(\C)$ record---in a precise sense---the Adams--Novikov spectral sequence for the sphere spectrum, including all differentials and extensions.  The ultimate expression of this connection is the discovery that the \emph{cellular} subcategory of $p$-complete $\C$-motivic spectra is equivalent to a category of \emph{synthetic} Adams--Novikov spectral sequences, and in particular has a purely homotopy theoretic description without reference to algebraic geometry.

%% The relationship between $p$-complete, cellular $\C$-motivic spectra and Adams--Novikov spectral sequences is mediated by a class $\tau:\Ss^1 \to (\G_m)^{\wedge}_p$.  The cofiber of the $\tau$ map (often denoted $C\tau$) is an $\mathbb{E}_\infty$-algebra, with the property that cellular $C\tau$-modules record the purely algebraic derived category of $p$-complete $\BP_*\BP$-comodules.  On the other hand, the full subcategory of $\tau$-invertible, $p$-complete motivic $\C$-spectra is equivalent to the classical category of $p$-complete spectra.  One summarizes the situation by saying that $\SH(\C)^{cell,p}$ is a $\tau$-deformation, with generic fiber the category of $p$-complete spectra and special fiber the derived category of $p$-complete $\BP_*\BP$-comodules.

The first evidence that a large sector of the motivic category contains only topological information appears in the work of Thomason on algebraic K-theory and \'etale cohomology \cite{Thomason}. As a consequence of this work, Thomason proves that for well-behaved $\C$-schemes $X$, there is an equivalence $\mathrm{K}^{\mathrm{alg}} (X) [\beta^{-1}]^{\wedge} _{\ell} \simeq \mathrm{KU} (X(\C))^{\wedge} _{\ell}$, where $\beta$ is a certain Bott element. %% There is a precise sense in which inverting $\beta$ is equivalent to inverting $\tau$.
These ideas have been refined over the years, culminating in the following theorem.

\begin{thm}[\cite{VoeOpenProbs, DI, LevineComparison, LevineSlicemU, Ctau, ctauactalol, Pstragowski, cmmf}]\label{thm:C-main}\

  In $\SHC$
  there is a map $\tau : \Ss^1 \to (\G_m)_p$
  which enjoys the following properties:
  \begin{enumerate}
  \item[(1.\hspace{1em}]\hspace{-1.1em}$\mathrm{Generically\ Topological})$ The full subcategory of $\tau$-local objects in $\SHC^{cell}_p$ is equivalent to $\Sp_p$.
  \item[(2.\hspace{1em}]\hspace{-1.1em}$\mathrm{Algebraic\ Degeneration})$ The cofiber of $\tau$ (often denoted $C\tau$) is a commutative algebra, and the category of dualizable modules over $C\tau$ is equivalent to the category of $p$-completions of dualizable objects in $\QCoh(\mathcal{M}_{\mathrm{fg}})$, where $\Mfg$ is the moduli stack of formal groups.
  \item[(3.\hspace{1em}]\hspace{-1.1em}$\mathrm{Galois\ Reconstruction})$ There is a purely topological construction of $\SHC_p^{cell}$.
  \end{enumerate}
  
  We summarize this situation by saying that $\SH(\C)_p^{cell}$ is a $1$-parameter deformation of $p$-complete spectra with parameter $\tau$ and a purely algebraic special fiber.
\end{thm}

Although it does not appear in the statement, understanding the algebraic cobordism spectrum, $\MGL$, is key to proving this theorem. 
The importance of $\MGL$ to the study of cellular motivic spectra goes back to Voevodsky's work on the effective slice filtration \cite{VoeOpenProbs}.
In that work, Voevosky conjectured that the effective slice filtration of the unit over an algebraically closed field 
may be described in terms of the Adams-Novikov spectral sequence. 
%% may be described in terms of the $\mathrm{E}_2$-page of the Adams-Novikov spectral sequence and further suggested a relationship between the effective slice spectral sequence of the unit over an algebraically closed field and the Adams-Novikov spectral sequence.
Levine later proved Voevodsky's conjectures \cite{LevineComparison, LevineSlicemU}.
Voevodsky's notion of ``rigid homotopy groups'' also foreshadowed \emph{algebraic degeneration}.
In this framework he and Rezk predicted that the rigid Adams spectral sequence is equivalent to the algebraic Novikov spectral sequence \cite[p. 20-21]{VoeOpenProbs}, a result later proven in different language by Gheorghe--Wang--Xu \cite[Theorem 1.17]{ctauactalol}.
The fact that the category of $p$-complete $\C$-motivic spectra is \emph{generically topological} is strongly suggested by results of Dugger--Isaksen \cite[Section 2.6]{DI}, though to the best of our knowledge the full result did not appear in print until \cite{Pstragowski}.\todo{I can't find an earlier one at least, but that could easily just be me. Piotr does beat Tom's etale rigidity paper.}
The commutative algebra structure on $C\tau$ was first constructed by Gheorghe \cite{Ctau},
and the category of modules over $C\tau$ was then identified by Gheorghe--Wang--Xu in \cite{ctauactalol}.
Finally, % \emph{Galois reconstruction} was accomplished independently by 
Pstragowski \cite{Pstragowski} and Gheorghe--Isaksen--Krause--Ricka \cite{cmmf}
independently provided two different \emph{Galois reconstructions} of cellular $p$-complete $\C$-motivic spectra. %% moreover, \cite{Pstragowski} reproves (2) independently of \cite{Ctau} and \cite{ctauactalol}.

%% In SECTION, we show that there is a precise dictionary between Voevodsky's language of the effective slice filtration and rigid homotopy groups, and the language used in \Cref{thm:C-main}.

With the case of $\C$ resolved we raise the following natural question:

\begin{qst}\label{qst:main}
  To what extent can this theorem be extended to a general field $k$?
\end{qst}

As stated, this question is too imprecise to admit a definitive answer, so we begin by refining two points of ambiguity: the appropriate subcategory of $\SHk$ one should consider and the meaning of purely topological. In our study of this question we have found that the appropriate subcategory is the category of Artin--Tate motivic spectra (defined below). Notably, over $\R$, if one restricts to the further subcategory of cellular objects, it becomes significantly more difficult to construct a topological model.

\begin{dfn}
  The category of Artin--Tate motivic spectra over $k$ is the smallest stable full subcategory
  $\SHkat \subset \SHk$,
  closed under tensor products and colimits, that contains the motives of finite \'etale $k$-algebras, 
  $\P_k^1$ and $(\P_k^1)^{-1}$.
\end{dfn}

The phrase `purely topological' has a double meaning.
On the one hand it refers to the input to the construction;
it should not depend directly on the arithmetic of $k$, instead using only the absolute Galois group, $G$, together with the character $G \to \widehat{\Z}^\times$ induced by the maximal cyclotomic extension of $k$ \footnote{More specifically, the $p$-complete category should only depend on the $\Z_p^\times$ component of this map.}.
On the other hand it asks for a construction which uses only the ``standard machinery of homotopy theory.''

The authors are not the first to take up questions of this nature.
%% Indeed, beyond the work by many authors in the case $k = \C$ discussed above,
Positselski has studied the question of Galois reconstruction for the category $\DM(k;\F_\ell)^{\at}$ of mixed Artin-Tate motives with mod $\ell$ coefficients \cite{DMRecon,GalKosz}. With a few exceptions, he has shown that when $k$ is a finite, local or global field, then $\DM(k,\F_\ell)^{\at}$ may be viewed as a derived category of filtered discrete $G$-modules with restricted sub-quotients. %the restriction being given in terms of the cyclotomic character.

In a different direction, work of Bachmann, Elmanto and {\O}stv{\ae}r shows that, up to a completion, $\SH(S)$ is generically \'etale for a wide range of schemes $S$ \cite{TomEtaleI, TomEtaleII}. In particular, Bachmann--Elmanto--{\O}stv{\ae}r show that, after a suitable completion, \'etale localization corresponds to inverting $\tau$. Note that $\tau$ may not exist in the homotopy of the completed unit, so care must be taken to interpret this statement.
Furthermore, Bachmann showed that, again up to completion, the \'etale motivic category is equivalent to the category of hypercomplete sheaves of spectra on the small \'etale site.
Specializing to the case $S = \Spec k$, we find that, up to completion, the category of $\tau$-local objects in $\SH(k)$ admits a description in terms of Borel $G$-equivariant spectra.

Previously, work of Behrens and Shah \cite{BehrensShah} had taken up the question of when a suitable map $\tau$ exists over $\R$. 
Although there is no map $\tau:\Ss^1 \to (\G_m)^{\wedge}_2$ in $\SH(\R)$,
%% though such a class does exist after tensoring with the spectrum $\HFt$ representing mod $2$ motivic cohomology.
they prove that $\tau$ exists whenever a different class $\rho:\Ss^0 \to \G_m$ has been killed.
If only $\rho^2$ is killed, but not $\rho$ itself, then $\tau$ does not necessarily exist, but $\tau^2$ does.
Continuing in this way, they make sense of inverting $\tau$ in any $\rho$-complete situation.
%% but the resulting $\tau$-inverted objects are more closely related to Borel $C_2$-spectra than genuine $C_2$-spectra.
%% Furthermore, there is no way to construct $C\tau$ without first killing $\rho$. This makes it impossible to describe cellular $\R$-motivic spectra as a $\tau$-deformation.

%% \todo{What's below needs to be joined up with what's above.}

%% In recent years the category $\SHR$ of $\R$-motivic spectra has similarly been found to shed light on the genuine $C_2$-equivariant stable homotopy category.  However, the exact relationship between $\R$-motivic stems and any sort of equivariant Adams--Novikov spectral sequence has remained mysterious, and before the present work there has been no purely homotopy theoretic construction of the cellular subcategory of $\SH(\R)$.

In this paper we resolve \Cref{qst:main} in the case $k=\R$.
This begins with the observation that if one does not insist that the target of $\tau$ must be a completion of $\G_m$, then a suitable replacement can easily be constructed.
%% The first 
%% The main moral of this paper is that all of these subtleties disappear so long as one is willing to consider grading $\R$-motivic homotopy groups not just by the Picard elements $\Ss^1$ and $\mathbb{G}_m$, but also by a third Picard element $\Ss^{\C}$.  Here, $\Ss^{\C}$ is the suspension spectrum of the unreduced suspension of $\Spec(\C)$.

\begin{thm}\label{thm:main}
  In $\SHRat$
  there is an invertible object $Q$ and a map $\ta : \Ss^1 \to Q_2$
  which enjoys the following properties:%\footnote{As we will explain in our conventions below, we use $\SHRat$ to denote the category of modules over the $2$-completion of the unit in $\SHR^{\at}$.}
  \begin{enumerate}
  \item[(GT)] The full subcategory of $\ta$-local objects in $\SHR^{\at}_{i2}$ is equivalent to $\Sp_{C_2,i2}$.
  \item[(AD)] The cofiber of $\ta$ (denoted $C\ta$) is a commutative algebra, and the category of dualizable modules over $C\ta$ is equivalent to the category of dualizable objects in the derived category of Mackey-functor $\MU_*\MU$-comodules \footnote{See \Cref{sec:cta} for a precise definition of this category.}.
  \item[(GR)] There is a purely topological construction of $\SHRat$.
    In particular, we construct a commutative algebra $R_\bullet$ in filtered, $C_2$-equivariant spectra
    such that the category of filtered modules over $R_\bullet$ is equivalent to $\SHRat$.
    The commutative algebra $R_\bullet$ is the even slice--d\'ecalage of the $\MUR$-Adams tower for the sphere \footnote{See \Cref{sec:top} for precise definitions.}.    
  \end{enumerate}
  
  We summarize this by saying that $\SHR^{\at}_{i2}$ is a $1$-parameter deformation of $C_2$-equivariant stable homotopy theory with parameter $\ta$ and purely algebraic special fiber.
\end{thm}

The romanization of the Devanagari letter $\ta$ is `ta' \footnote{However, the reader who is not familiar with the Devanagari alphabet is advised that the pronunciation of $\ta$ in the majority of South Asian languages using Devanagari is given in IPA by /\textsubbridge{t}\textschwa/, which is closer to `tuh' in English than `ta.'}  and
our choice of this symbol will be explained in \Cref{rmk:ta-explanation}. As in the case over $\C$, we are also able to give an explicit formula for the homotopy groups of $C\ta$, though since we have not yet set up the appropriate notion of homotopy groups we will defer an explicit statement until later.
%% Note that the $p$-completions present in the theorem statement are necessary.
The subscript $i2$ appearing in the theorem statement refers to the category of modules over the $2$-completion of the unit. 
The reader who is worried by this departure from the norm may wish to read the section on notations and conventions before proceeding.
%% about the handling of $p$-completion may wish to read \Cref{app:completion} which sets up our general approach to completion issues before proceeding.

\ 

With the resolution of \Cref{qst:main} in the case of $\R$, the authors are hopeful that we might see this question resolved in its entirety in the near future.

%% \begin{cnstr}
%% There is a canonical class $\ta: \Ss^{\C} \to (\G_m)^{\wedge}_2$.  We will build $\ta$ in \Cref{sec:tau} via the explicit algebraic geometry of the circle $Q=\Spec \R[x,y]/(x^2+y^2-1)$.  A key point is that $\Sigma^{\infty} Q$ is itself a Picard element in $\SHgc$, satisfying a relation $\Ss^{\C} \otimes \Sigma^{\infty} Q \simeq \Ss^1 \otimes \G_m$.
%% \end{cnstr

%% \begin{rmk}
%%   The romanization of the Devanagari letter $\ta$ is `ta.' \footnote{However, the reader who does not know Hindi or another language which utilizes the the Devanagari alphabet should be advised that its pronunciation is closer to `tuh.' To be even more correct, the `t' sound should be dental and unaspirated.}  As we will explain in \Cref{sec:homology}, our choice of this character arises from the fact that there is a class $u:\Ss^1 \otimes (\Ss^{\C})^{-1} \to \HFt$ and a product relation $\ta u = \tau \in \pi_{***}^{\R}(\HFt).$    In other words, ta $\cdot$ u is tau.
	
%% It follows that inverting $\tau$ is the same as inverting both $\ta$ and $u$.  We have shown that inverting $\ta$ is equivalent to Betti realization, while inverting $u$ Borel completes genuine $C_2$ spectra.


\subsection{Examples of Artin--Tate motivic spectra}\ 

In this subsection, which is preparatory to all later material, we give a more comprehensive introduction to the category of Artin--Tate motivic spectra over a field. This begins with introducing the smaller subcategories of Artin and Tate objects, which together generate the category of Artin--Tate objects. The bulk of the subsection is spent building a collection of important examples of Artin--Tate motivic spectra. Over $\R$ a special role is played by the invertible objects and so we focus specific attention there. We close the section by discussing homotopy groups.

\begin{dfn}
  The categories of Artin, Tate and Artin--Tate motivic spectra over $k$ are the stable, full subcategories of $\SHk$, closed under tensor products and colimits, with the following generators:
  \begin{itemize}
  \item The category of Artin motivic spectra, $\SHka$, is generated by the motives of finite \'etale $k$-algebras.
  \item The category of Tate motivic spectra, $\SHkt$, is generated by the motives of $\P_k^1$ and $(\P_k^1)^{-1}$ \footnote{This category is sometimes referred to as the cellular category following \cite{DICell}.}.
  \item The category of Artin--Tate motivic spectra, $\SHkat$, is generated by the motives of finite \'etale $k$-algebras together with $\P_k^1$ and $(\P_k^1)^{-1}$.
  \end{itemize}
  Since these subcategories are closed under tensor products and colimits, they each inherit the structure of a stable, presentably symmetric monoidal category from $\SHk$. \footnote{In the sequel, we will state without remark results that are only known to be true after the characteristic of $k$ is inverted. The reader is invited to restrict themself to $k$ of characteristic zero if they prefer.}
\end{dfn}

We now turn to examples of objects in each of these categories, starting with the trivial and heading towards the non-trival.

\begin{exm} \label{exm:Sp-unit}
  Every stable, presentably symmetric monoidal category admits a unique symmetric monoidal functor from the category of spectra \cite[Corollary 4.8.2.19]{HA} \footnote{Another way of saying this is that $\Sp$ is the initial object of $\mathrm{CAlg}(Pr^{L,\mathrm{stab}})$.}.
  This gives us objects $X \otimes \o_k$ for every spectrum $X$.
  Of particular interest are the integer simplicial suspensions of the unit, $\Ss^{n} \otimes \o_k$.
\end{exm}

Inductively applying the homotopy purity theorem \cite[Theorem 3.2.23]{MorelVoevodsky} to the decomposition $\P^n _k = \mathbb{A}^{n} _k \coprod \P^{n-1} _k$, we learn that:

\begin{exm}[{\cite[Example 2.12]{DICell}}]
  For each $n$, the motivic spectrum associated to $\P^n_k$ is Tate.
\end{exm}

More generally, using a Bialynicki-Birula decomposition, Wendt shows that any smooth projective variety which admits a $\G_m$-action whose fixed points are discrete and rational is Tate \cite{Wendt} \footnote{In the case where the fixed points are not rational, the authors wonder what conditions are necessary to guarantees the motive is Artin--Tate.}.

\begin{exm}[{\cite[Theorem 6.4]{DICell}, \cite[Proposition 8.1]{HM}, \cite[Proposition 8.12]{etaPeriodic}}]
  Each of the following commutative algebras is Tate:
  \[ \MGL,\quad \kq,\quad \kgl,\quad \mathrm{M}\Z,\quad \mathrm{M}\F_p. \]
\end{exm}

\begin{exm} \label{exm:galois-corr}
  Essentially by definition, the Galois correspondence provides functors
  \begin{center}
    \begin{tikzcd}
      {\left\{ \footnotesize\parbox{2.25cm}{finite, continuous Gal($\overline{k}/k$)-sets} \right\}} \ar[r] &
      {\left\{ \footnotesize\parbox{1.5cm}{finite, etale $k$-algebras} \right\}^{op}} \ar[r, hook] \ar[d] &
      \mathrm{Sm}_k \ar[d] \\
      & \SHka \ar[r, hook] &
      \SHk,
    \end{tikzcd}
  \end{center}
  which give our first examples of Artin objects.
\end{exm}

Since the functor in the example above sends disjoint unions to sums there is no loss in generality if we restrict to transitive Gal($\overline{k}/k$)-sets (field extensions). In the case of $\R$ this gives only two objects: $\Spec(\R)$, which we will temporarily denote $\o_\R$ since is it the monoidal unit of the category, and $\Spec(\C)$. Using the natural map we get a cofiber sequence
\begin{align} \Spec(\C) \to \o_\R \xrightarrow{a} \Ss^{\C}. \label{eq:pic-C} \end{align}
On the level of $\C$-points (with Galois action), this cofiber sequence gives the representation sphere $\Ss^{\sigma}$.




Given a quadric in the plane, $V$, we can take its closure in $\P_k^2$ to obtain $\overline{V}$ which is a form of $\P_k^1$. If we let $x_1,\dots,x_n$ denote the points at infinity, then by the homotopy purity theorem \cite[Theorem 3.2.23]{MorelVoevodsky} we have a cofiber sequence
\[ V \to \overline{V} \to \bigoplus_{i=1,\dots,n} \P_{k(x_i)}^1. \]
Under the assumption that $V$ admits a rational point, $\overline{V}$ is just $\P^1_k$, and so we obtain our first non-trivial example of a motive which is neither Artin nor Tate, but is Artin--Tate:

\begin{exm}
  The motive of an affine quadric which admits a rational point is Artin--Tate.
\end{exm}

Over $\R$ there is a particular affine quadric of interest to us.
It is the object $Q$ which appeared in the statement of \Cref{thm:main}.

\begin{exm}
  Let $Q \coloneqq \{x^2+y^2=1\} \subset \AA^2 _{\R}$, which we shall call the algebraic circle.
  %% From the previous example we know that $Q$ is Artin--Tate.
  $Q$ has the additional property that it is a form of $\G_m$ \footnote{By Cartier duality, forms of $\G_m$ are classified by rank 1 lattices with Galois action. $Q$ corresponds to the unique nontrival action.}.
  Its group scheme structure comes from the usual rule for multiplication of complex numbers,
  namely
  \[ (x_1,y_1) \cdot (x_2,y_2) = (x_1 x_2 - y_1 y_2,\ x_1 y_2 + y_1 x_2). \]
  %% Furthermore, it is a form of $\G_m$---there is an isomorphism $(Q)_\C \cong (\G_m)_\C$ of commutative group schemes over $\C$.

  In \Cref{sec:ta} we will construct the maps $\ta: \Ss^1 \to Q_p$ using our explicit understanding of this group scheme.
\end{exm}

\begin{prop}[{\cite[Proposition 1.1]{HuPicard}}] \label{prop:pic-Q}
  There is an equivalence $Q \otimes \Ss^\C \simeq \P^1 _\R$.
  In particular, both $\Ss^\C$ and $Q$ are invertible.
\end{prop}

At this point we are ready to define the appropriate notion of homotopy groups for studying $\SHR^{\at}$.
In the case of $\SHC^{\at}$, the category has a bi-graded family of compact invertible generators given by the spheres $\Ss^{p - 2w} \otimes (\P_\C^1)^{\otimes w}$. Therefore, it is typical to study objects at the level of their bigraded homotopy groups, given by
\[ \pi_{p,w}^\C(X) \coloneqq \pi_0\Map ( \Ss^{p-2w} \otimes (\P_\C^1)^{\otimes w}, X ). \]
Similarly, $\SHRt$ has a bi-graded family of compact invertible generators given by the spheres $\Ss^{p-2w} \otimes (\P_\R^1)^{\otimes w}$, and it is typical to study objects through their bi-graded homotopy groups.
Using \Cref{prop:pic-Q} and \Cref{eq:pic-C}, the category $\SHR^{\at}$ has a tri-graded family of compact invertible generators given by the spheres $\Ss^{p-w} \otimes (\Ss^\C)^{\otimes q-w} \otimes (\P_\R^1)^{\otimes w}$, and we will study most objects through their tri-graded homotopy groups.

\begin{ntn}
  We endow the categories $\Sp$, $\Sp_{C_2}$, $\SHC$ and $\SHR$ with Picard gradings by spheres as follows:
  \begin{center}
    \renewcommand{\arraystretch}{1.1}
    \begin{tabular}{|l|l|l|} \hline
      Category & Picard & Spheres \\\hline\hline
      $\Sp$ & $\Z$ & $\Ss^p \simeq (S^1)^{\otimes p}$ \\\hline
      $\Sp_{C_2}$ & $\Z \times \Z$ & $\Ss^{p+q\sigma} \simeq (S^1)^{\otimes p} \otimes (S^\sigma)^{\otimes q}$ \\\hline
      $\SHC$ & $\Z \times \Z$ & $\Ss^{p,w} \simeq (S^1)^{\otimes p-2w} \otimes (\P_\C^1)^{\otimes w}$ \\\hline
      $\SHR$ & $\Z \times \Z \times \Z$ & $\Ss^{p,q,w} \simeq (S^1)^{\otimes p-w} \otimes (S^\C)^{\otimes q-w} \otimes (\P_\R^1)^{\otimes w}$ \\\hline
    \end{tabular}
    \renewcommand{\arraystretch}{1.0}
  \end{center}
  
  For any $p,q,w \in \Z$, and any $\R$-motivic spectrum $X$, we let $\pi^{\R}_{p,q,w} X$ denote the group of homotopy classes of maps,
  \[ (S^1)^{\otimes p-w} \otimes (S^\C)^{\otimes q-w} \otimes (\P_\R^1)^{\otimes w} \to X. \]
\end{ntn}

\begin{exm}
  For convenience we record what the invertible objects considered in this section look like under this new notation:
  \begin{center}
    \begin{tabular}{ccc}   
      $\Ss^{0,0,0} \cong \o_\R$ \qquad\qquad\qquad &
      $\Ss^{1,0,0} \cong S^1$ \qquad\qquad\qquad & 
      $\Ss^{0,1,0} \cong S^{\C}$ \\
      $\Ss^{0,1,1} \cong \G_m$ \qquad\qquad\qquad &
      $\Ss^{1,0,1} \cong Q$ \qquad\qquad\qquad &
      $\Ss^{1,1,1} \cong \P_\R^1$
    \end{tabular}
  \end{center}
  In the case of $\Ss^{0,0,0}$ we will often drop the indices for brevity, writing only $\Ss$.
  The two maps between Picard elements considered thus far, $\ta$ and $a$, become
  \[ \ta \in \pi_{0,0,-1}^\R \Ss_p \qquad\quad \text{ and } \qquad\quad a \in \pi_{0,-1,0}^\R \Ss. \]
  %% In the case when $X$ is the unit we will often omit it from the notation for brevity.
\end{exm}

\begin{rmk}
  The tri-graded spheres in $\SHR^{\at}$ which are Tate are precisely the spheres of the form $\Ss^{p,q,q}$.
  For this reason the Tate category only sees a `slice' of the total information present in the tri-graded homotopy groups.

  The tri-graded spheres in $\SHR^{\at}$ which are Artin are precisely the spheres of the form $\Ss^{p,q,0}$.
  The Artin category similarly only sees a `slice' of the total homotopy groups.

  %% Traditionally, $\SHR$ is viewed as bigraded by the spheres $\Ss^{p,w} \simeq (S^1)^{\otimes p-w} \otimes (\G_m)^{w}$. The comparison with our trigraded spheres is given by the equivalence $\Ss^{p,w} \simeq \Ss^{p-w,w,w}$, and from this we see that the bigraded sheres capture precisely those trigraded spheres $\Ss^{p,q,w}$ satisfying $q=w$.
\end{rmk}

%% as a  defined by [people]. We regard this category as the correct analog of the category over cellular objects over an algebraically closed field and propose that it may be possible to extend the full package of results known in the cellular case over algebraically closed fields to the Artin-Tate category over general fields. 


%% the the Picard elements $\Ss^1$, $\G_m$, $\Ss^{\C}$ and their inverses.

%% Let $\Ss_2^{\wedge}$ denote the $2$-completed unit in $\SHR$, and consider the category $\Ss_2^{\wedge}$-mod of $\Ss_2^{\wedge}$-module spectra.  We similarly define $\SHgc_2$ to be the smallest full subcategory of $\Ss_2^{\wedge}$-mod, closed under colimits and tensor products over $\Ss^{\wedge}_{2}$, containing the Picard elements $(\Ss^1)^{\wedge}_2, (\G_m)_2^{\wedge}$, $(\Ss^{\C})_2^{\wedge}$, and their inverses.

%% Both $\SHgc$ and $\SHgc_2$ are symmetric monoidal stable $\infty$-categories, and by CITE the $2$-completions of these categories are equivalent.  Note however that not every object in $\SHgc_2$ is $2$-complete.
%% \end{dfn}


%% \begin{cnstr}
%% There is a canonical class $\ta: \Ss^{\C} \to (\G_m)^{\wedge}_2$.  We will build $\ta$ in \Cref{sec:tau} via the explicit algebraic geometry of the circle $Q=\Spec \R[x,y]/(x^2+y^2-1)$.  A key point is that $\Sigma^{\infty} Q$ is itself a Picard element in $\SHgc$, satisfying a relation $\Ss^{\C} \otimes \Sigma^{\infty} Q \simeq \Ss^1 \otimes \G_m$.
%% \end{cnstr}


%% In this section, we introduce some exotic Picard elements in $\SHR$, use them to define trigraded homotopy groups, and define some maps between them, in particular the element $\ta \in \pi^{\R} _{0,0,-1} \Ss_p$, which serves as a Picard-graded analogue of the class $\tau \in \pi^{\C} _{0,-1} \Ss_p$.

%% \todo{Give a general definition of things}
%% \todo{Conjecture: every pic element is already in this category}

%% We discuss Picard gradings in $\Sp$, $\Sp_{C_2}$, $\SHC$ and $\SHR$. We begin by introducing the invertible objects that will generate our Picard gradings.

%% \begin{rec}
%%   Given an category $\CC$ admitting colimits, we let $S^1$ denote the simplicial circle, i.e. the colimit of the constant diagram on $S^1$ at the terminal object.
%%   In $\mathrm{Spaces}_{C_2}$, let $S^{\sigma}$ denote the unreduced suspension of $C_2$. Similarly, in $\MR$, we let $S^{\C}$ denote the unreduced suspension of $\Spec \C$.
%%   Finally, in $\MR$ and $\MC$, we let $\G_m$ denote the multiplicative group scheme.
%%   By abuse of notation, we will use the same notation for the suspension spectra of $S^1$, $S^{\sigma}$, $S^\C$ and $\G_m$.

%%   Then $S^1$ is invertible in any stable category.
%%   It is well-known that $S^{\sigma}$ is invertible in $\Sp_{C_2}$, and that $\G_m$ is invertible in $\SHC$ and $\SHR$.
%%   It is a result of Hu that $S^\C$ is invertible in $\SHR$ \cite[Proposition 1.1]{HuPicard}.
%% \end{rec}


%
%\begin{rec}
  %Let $\sigma$ denote the real sign representation of $C_2$, and let $S^\sigma$ denote the one-point compactification of $\sigma$. There is also a cofiber sequence of pointed $C_2$-spaces $(C_2)_+ \to S^0 \to S^{\sigma}$.
%
  %The stabilized sphere $\Ss^{\sigma}$ is invertible, and one defines $\Ss^{p+q\sigma} = (\Ss^1)^{\otimes p} \otimes (\Ss^\sigma)^{\otimes q}$. This endows the category of $C_2$-spectra with an $RO(C_2)$-grading, which may also be viewed as a bigrading.
%\end{rec}
%
%\begin{rec}
  %Let $S^1 \in \MC$ denote the simplicial circle and let $\G_m \in \MC$ be the motivic space associated to the multiplicative group. Let $\Ss^{1,0} \coloneqq \Sigma^\infty S^1$ and $\Ss^{1,1} \coloneqq \Sigma^\infty \G_m$. We then define $\Ss^{a,b} = (\Ss^{1,0})^{\otimes a-b} \otimes (\Ss^{1,1})^{\otimes b}$. This endows $\SHC$ with a bigrading.
%\end{rec}


%% \begin{rmk}
%%   Traditionally, $\SHR$ is viewed as bigraded by the spheres $\Ss^{p,w} \simeq (S^1)^{\otimes p-w} \otimes (\G_m)^{w}$. The comparison with our trigraded spheres is given by the equivalence $\Ss^{p,w} \simeq \Ss^{p-w,w,w}$, and from this we see that the bigraded sheres capture precisely those trigraded spheres $\Ss^{p,q,w}$ satisfying $q=w$.
%% \end{rmk}

%% There is one more explicit Picard element in $\SHR$ that shall be of interested to us.

%% \begin{prop}
%%   Stably, there is an equivalence $Q \simeq \Ss^{1,0,1}$. In particular, $Q$ is invertible.
%% \end{prop}

%% \begin{proof}
%%   Hu has shown that $Q \otimes \Ss^\C \simeq \P^1 _\R$ \cite[Proposition 1.1]{HuPicard}. On the other hand, it is standard that $\G_m \otimes \Ss^1 \simeq \P^1 _\R$, from which it follows that $\P^1 _\R \simeq \Ss^{1,1,1}$, which implies the desired result.
%% \end{proof}





%Traditionally, $\SHR$ is also studied via its bigraded homotopy groups. In this paper, we will endow $\SHR$ further with a trigrading.
%As our first step in introducing this trigrading, we recall several invertible $\R$-motivic spectra.
%
%\begin{dfn}
  %We define several $\R$-motivic spaces as follows:
  %\begin{itemize}
    %\item $S^1$ is the simplicial circle.
    %\item $\G_m$ is the multiplicative group.
    %\item $\P^1$ is the projective line.
    %\item $S^\C$ is the unreduced suspension of $\Spec \C$.
    %\item $Q$ is the affine quadric $\{x^2 + y^2 = 1\} \subseteq \AA^2$, which is an abelian group scheme under complex multiplication: $(x_1,y_1) \cdot (x_2,y_2) = (x_1 x_2 - y_1 y_2, x_1 y_2 + y_1 x_2)$.
  %\end{itemize}
  %We will abuse notation by using the same notation to denote the associated $\R$-motivic suspension spectra.
%\end{dfn}
%
%\begin{rmk}
  %Note that $Q$ is a form of $\G_m$: there is an isomorphism of group schemes $Q_\C \cong (\G_m)_\C$. Moreover, there is an equivalence of $\C$-motivic spaces $(S^\CC)_\CC \simeq (S^1)_\CC$.
%\end{rmk}
%
%\begin{prop}
  %The $\R$-motivic spectra $S^1$,$\G_m$, $\P^1$, $S^\C$ and $Q$ are all invertible. Moreover, there are relations $S^1 \otimes \G_m \simeq S^\C \otimes Q \simeq \P^1$.
%\end{prop}
%
%\begin{proof}
  %It is classical that $S^1$, $\G_m$ and $\P^1$ are invertible, as is the relation $S^1 \otimes \G_m \simeq \P^1$. The relation $S^\C \otimes Q \simeq \P^1$ is \cite[Proposition 1.1]{HuPicard} and implies that $S^\C$ and $Q$ are invertible.
%\end{proof}
%
%The $\R$-motivic spectra $S^1, \G_m$ and $\P^1$ are inverible by the definition of $\SHR$, and there is the classical relation $S^1 \otimes \G_m \simeq \P^1$ among them. In \cite{HuPicard}, it is shown that $S^\C \otimes Q \simeq \P^1$, showing that $S^\C$ and $Q$ are also invertible. After base change to $\C$, there are equivalences $S^{\C} _{\C} \simeq S^1 _\C$ and $Q_{\C} \cong (\G_m)_\C$, the latter originating from an isomorphism of group schemes.
%
%Using the above Picard elements and relations, we are able to introduce a trigrading on the $\R$-motivic stable homotopy category, extending the usual bigrading.
%
%\begin{ntn}
  %We define trigraded spheres as follows. First, we define:
  %\begin{itemize}
    %\item $\Ss^{1,0,0} \coloneqq S^1$.
    %\item $\Ss^{0,1,0} \coloneqq S^{\C}$.
    %\item $\Ss^{0,1,1} \coloneqq \G_m$.
    %\item $\Ss^{1,0,1} \coloneqq Q$.
    %\item $\Ss^{1,1,1} \coloneqq \P^1$.
  %\end{itemize}
  %In general, we define $\Ss^{p,q,w}$ for $(p,q,w) \in \Z^{\times 3}$ by demanding that $\Ss^{p_1,q_1,w_1} \otimes \Ss^{p_2,q_2,w_2} \cong \Ss^{p_1+p_2,q_1+q_2,w_1+w_2}$. This is possible because of the equivalences $S^1 \otimes \G_m \simeq S^\C \otimes Q \simeq \P^1$.
%\end{ntn}
%
%\begin{rmk}
  %Recall the usual motivic bigrading, where $\Ss^{1,0} \coloneqq S^1$ and $\Ss^{1,1} \coloneqq \G_m$. Then the bigraded sphere $\Ss^{a,b}$ is equivalent to the trigraded sphere $\Ss^{a-b,b,b}$. A trigraded sphere $\Ss^{p,q,w}$ is equivalent to a bigraded sphere precisely when $q=w$.
%\end{rmk}
%
%Under this trigrading, the first two gradings keep track of the underlying $RO(C_2)$-grading, and the third grading is the motivic weight.
%
%\begin{prop}
  %Under the Betti realization functor $\b : \SHR \to \Sp_{C_2}$, we have $\b(\Ss^{a,b,c}) \simeq \Ss^{a+b\sigma}$.
  %Under the base change functor $(-)_\C : \SHR \to \SHC$, we have $(\Ss^{a,b,c})_\C \cong \Ss^{a+b,c}$. \todo{What's a better way to denote this?}
%
  %As a consequence, we find that $\Ss^{p_1,q_1,w_1} \simeq \Ss^{p_2,q_2,w_2}$ implies that $(p_1,q_1,w_1) = (p_2,q_2,w_2)$.
%\end{prop}

%% The categories $\Sp$ and $\Sp_{C_2}$ are generated under colimits by the Picard spheres generated above. This is not true of $\SHC$ and $\SHR$, so if we want to study these categories via the Picard graded homotopy groups indexed by the spheres considered above, we would do well to restrict ourselves to the full subcategories generated by these spheres.



\subsection{Comparison functors}\ 

The two simplest ways to interrogate a category are to study specific objects % of note inside the category
and to study the network of functors which relate it to other categories.
While the previous subsection provided preperatory background on specific objects, this subsection sets up the suite of functors which we will use to produce and study more general objects.
In the specific case of $\R$ this means studying the various ways we can move between $\SHR$, $\SHC$, $\Sp_{C_2}$ and $\Sp$.

%% Though parts of this subsection may appear unmotivated, it provides the background necessary to begin seriously discussing the structure of $\SHRat$ in the following subsection.

%% We begin the subsection by reviewing the functors provided by the six functor formalism.
%% Next, we discuss Betti realization over $\C$ and $\R$.
%% One of the surprising features of these realization functors that doesn't extend to more general fields is that they admit a symmetric monoidal section.
%% DISCUSS NU
%% DISCUSS NORMS

\begin{rec} \label{rec:fields}
  Given a finite extension of fields $i : k \to \ell$, there are two pairs of adjunctions and one extra functor coming from the six functor formalism
  \[ i^* : \SHk \rightleftarrows \SH(\ell) : i_*, \qquad i_! : \SH(\ell) \rightleftarrows \SHk : i^! \quad \mathrm{ and } \quad i_{\sharp} : \SH(\ell) \to \SHk \]
  where $i^*$ is symmetric monoidal.
  Since $i$ is smooth, proper and unramified we have equivalences $i_! \simeq i_* \simeq i_{\sharp}$ and $i^! \simeq i^*$ \cite[Theorems 6.9 and 6.18]{hoyois}.
  These equivalences tell us that
  \begin{enumerate}
  \item $i^*$ and $i_*$ both commute with all limits and colimits.
  \item If $X$ is a smooth $\ell$-scheme, then $i_*X \simeq X$, where the second copy of $X$ is considered as a $k$-scheme.
  \end{enumerate}
  \todo{Restrictions ?}
\end{rec}

In \Cref{app:boilerplate}, as an example of the techniques showcased there, we show that for fields of characteristic zero
there is an equivalence of presentably symmetric monoidal categories
\[ \SH(\ell) \simeq \Mod ( \SHk ; \Spec(\ell) ), \]
and $i_*i^*(-) \simeq \Spec(\ell) \otimes -$.
This equivalence tells us that no new information enters the picture when we pass to a field extension.
Specializing to the case of $\C/\R$ the above equivalence becomes
\[ \SHC \simeq  \Mod ( \SHR ; \Spec(\C) ). \]

Using the descriptions of $i_*$ and $i_*i^*$ we can conclude that both $i^*$ and $i_*$ restrict to the full subcategory of Artin--Tate objects. Thus we obtain a similar description of the Artin--Tate category of a field extension:
\[ \SH(\ell)^{\at} \simeq \Mod ( \SHkat ; \Spec(\ell) ). \]
Specializing to the case of $\C/\R$ we will sometimes denote $\Spec(\C)$ by $Ca$ since it is (by the definition of $a$) the cofiber of the map $a: \Ss^{0,-1,0} \to \Ss^{0,0,0}$.
The above equivalence becomes
\[ \SHC^{\at} \simeq \Mod ( \SHR^{\at} ; Ca ). \]

At this point we turn to studying the case of $\R$ more closely.
Though some of the things we do after this point have obvious analogs in other cases,
many of our key maneuvers implicitly rely on the fact that $\R$ has a finite (and well-understood) absolute Galois group.
In particular, we will now assume that the reader is familiar with $C_2$-equivariant homotopy theory as in \cite{HHR}.
Equivariant homotopy theory first enters the picture through the Betti realization functors of \cite[Section 3.3]{MorelVoevodsky}.

\begin{rec}

  There is a commutative diagram of symmetric monoidal left adjoints:
  %% \begin{center}
  %%   \begin{tikzcd}[column sep=tiny, row sep=tiny]
  %%     \Ss^{p,q,w} \ar[rrrrrrr,mapsto] \ar[ddddddd,mapsto] & & & & & & & \Ss^{p+q,w} \ar[ddddddd,mapsto] \\
  %%     & \SHR \ar[rrrrr,"(-)_\C"] \ar[ddddd,"\b"] & & & & & \SHC \ar[ddddd,"\b"] & \\
  %%     & & & & & & & \\
  %%     & & & & & & & \\
  %%     & & & & & & & \\
  %%     & & & & & & & \\
  %%     & \Sp_{C_2} \ar[rrrrr,"\Phi^e"] & & & & & \Sp & \\
  %%     \Ss^{p+q\sigma} \ar[rrrrrrr,mapsto] & & & & & & & \Ss^{p+q}.
  %%   \end{tikzcd}
  %% \end{center}

  \begin{center}
    \begin{tikzcd}[column sep=tiny, row sep=tiny]
      
      \SHR \ar[rrrrrrr,"(-)_\C"] \ar[ddddddd,"\b"] & & & & & & & \SHC \ar[ddddddd,"\b"] \\
      & \Ss^{p,q,w} \ar[rrrrr,mapsto] \ar[ddddd,mapsto] & & & & & \Ss^{p+q,w} \ar[ddddd,mapsto] \\
      & & & & & & & \\
      & & & & & & & \\
      & & & & & & & \\
      & & & & & & & \\
      & \Ss^{p+q\sigma} \ar[rrrrr,mapsto] & & & & & \Ss^{p+q} \\
      \Sp_{C_2} \ar[rrrrrrr,"\Phi^e"] & & & & & & & \Sp.

    \end{tikzcd}
  \end{center}
  
  %% \begin{center}
  %%   \begin{tikzcd}
  %%   \SHR \ar[rrr,"(-)_\C"] \ar[ddd,"\b"] & & & \SHC \ar[ddd,"\b"] \\  
  %%   & \Ss^{p,q,w} \ar[r, mapsto] \ar[d, mapsto] & \Ss^{p+q,w} \ar[d, mapsto] \\
  %%   &  \Ss^{p+q\sigma} \ar[r, mapsto] &\Ss^{p+q} \\
  %%   \Sp_{C_2} \ar[rrr,"\Phi^e"] & & & \Sp     
  %%   \end{tikzcd}
  %% \end{center}

  In the above diagram, $(-)_\C$ is the base change functor, $\Phi^e$ is the underlying functor and $\b$ are the Betti realization functors, induced by the assigment $X \mapsto X(\C)$.
  The inner square describing these functors on Picard objects can be verified by considering $Q$, $\G_m$ and $S^1$ directly \footnote{The observation that spurred the authors to begin this project was that upon restricting to categories of Tate objects the induced square on Picard groups is not a pullback, but with Artin--Tate objects it is in fact a pullback.}.
\end{rec}

\begin{rmk}
  Since the Betti realization of $\Ss^{p,q,w}$ is $\Ss^{p+q\sigma}$, we think of $w$ as recording the motivic weight and $p,q$ as providing a copy of $RO(C_2)$ in each weight.
  %% We also see that $\abs{\ta} = (0,0,-1)$, so that $\ta$ shifts only the weight grading, whereas $\abs{\tau} = (1,-1,-1)$ also shifts the $RO(C_2)$-grading.
  The Tate spheres are those of the form $\Ss^{p,q,q}$, so in the bigraded world the number of $\sigma$'s in the $RO(C_2)$-grading must equal the motivic weight.
  This restriction blinds one to the existence of a weight shifting element $\ta$ in $RO(C_2)$ grading $0$, which is the true analog of the $\C$-motivic $\tau$.
\end{rmk}

Having discussed Betti realization, we now introduce the less well-known functors $c$ and $c_{\C/\R}$ which provide sections of Betti realization.

%% functors out of $\SHR$ we now discuss the various functors which allow us to produce objects in $\SHR$. This begins with the symmetric monoidal left adjoints $c$ and $c_{\C/\R}$ which provide sections of the Betti realization functors. Then, we discuss the more exotic operations given by the functor $\nu$ and the norm functors of [CITE], neither of which are exact. It is the functor $\nu$ which is most useful for producing motivic analogs of topological objects and whose construction is most delicate.

\begin{rec} \label{rec:c}
  There are symmetric monoidal left adjoints $c$ and $c_{\C/\R}$,
  which fit into the commutative diagram
  \begin{center}
    \begin{tikzcd}[sep=huge]
      \Ss^{p+q\sigma} \ar[r, mapsto] & \Ss^{p,q,0} \ar[r, mapsto] & \Ss^{p+q\sigma} \\
      \Sp_{C_2} \ar[rr, "\mathrm{id}_{\Sp_{C_2}}", bend left=20] \ar[r, "c_{\C/\R}"] \ar[d,"\Phi^e"] & \SHR \ar[r, "\b"] \ar[d,"(-)_\C"] & \Sp_{C_2} \ar[d,"\Phi^e"] \\
      \Sp \ar[rr, "\mathrm{id}_{\Sp}"', bend right=20] \ar[r, "c"] & \SHC \ar[r, "\b"] & \Sp \\
      \Ss^{p} \ar[r, mapsto] & \Ss^{p,0} \ar[r, mapsto] & \Ss^p .
    \end{tikzcd}
  \end{center}

  The functor $c$ is the unique symmetric monoidal left adjoint coming from the fact that $\Sp$ is the unit of $Pr^{L, \mathrm{stab}}$ (see \Cref{exm:Sp-unit}).
  The functor $c_{\C/\R}$ is constructed in \cite{HellerOrmsby}, by beginning with the functor
  \[ \left\{ \mathrm{finite} \ \ \  C_2\mathrm{-sets} \right\} \to \SHR \]
  from \Cref{exm:galois-corr} and then extending it to all of $\Sp_{C_2}$.
  %% which captures the $C_2$-Galois extension $\C/\R$, sending $C_2/e$ to $\Spec \C$ and $C_2/C_2$ to $\Spec \R$. NAME then shows that this extends to the functor $c_{\C/\R} : \Sp_{C_2} \to \SHR$.
\end{rec}

Since they are symmetric monoidal left adjoints, it is easy to see that the composites $\b \circ c : \Sp \to \Sp$ and $\b \circ c_{\C/\R} : \Sp_{C_2} \to \Sp_{C_2}$ are equivalent to the identity. In fact, upon restricting to the appropriate target category we uncover something more interesting:

\begin{thm}[{\cite{LevineComparison,HellerOrmsby,HellerOrmsbyII}}] \label{thm:c-ff}
  The symmetric monoidal functors $c$ and $c_{\C/\R}$ factor through the respective categories of Artin objects and provide equivalences,
  \[ c : \Sp \xrightarrow{\simeq} \SHCa \qquad \text{ and } \qquad c_{\C/\R} : \Sp_{C_2} \xrightarrow{\simeq} \SHRa. \]
\end{thm}

\begin{cor} \label{cor:betti-iso}
  The induced maps,
  \[ \pi_p\Ss \xrightarrow{c} \pi^\C _{p,0}\Ss \xrightarrow{\b} \pi_p\Ss  \qquad \text{and} \qquad \pi^{C_2}_{p+q\sigma}\Ss \xrightarrow{c_{\C/\R}} \pi^\R _{p,q,0}\Ss \xrightarrow{\b} \pi^{C_2} _{p+q\sigma}\Ss \]
  are isomorphisms for all $p$ and $q$.
\end{cor}

\begin{rmk}
  This corollary provides an identification of $\pi_{0,0,0}^{\R}\Ss$ with the Burnside ring of $C_2$.
  It is isomorphic to $\Z \oplus \Z$, with generators $1$ and $[C_2]$ and the relation $[C_2]^2 = 2[C_2]$.

  Morel's identification of $\pi_{n,n}\o_k$ in terms of Milnor--Witt K-theory provides another way of assigning names to elements of $\pi^\R_{0,0,0}\Ss$ \cite{Morelpizero}.
  In these terms $\pi_{0,0,0}^\R\Ss$ is generated by $1$ and $\eta [-1]$, subject to the relation $(\eta [-1])^2 = -2 \eta [-1]$. We denote by $\rho$ the element $-[-1]$.
  The translation between the two bases described above is given by $\eta \rho = 2 - [C_2]$.
\end{rmk}

%% It is easy to see from the definitions that the functors $c$ and $c_{\C/\R}$ land in the subcategories $\SHC^{cell}$ and $\SHgc$. Indeed, their essential images are the full subcategories generated under colimits by the weight $0$ spheres $\Ss^{p,0}$ and $\Ss^{p,q,0}$, respectively.

%% \begin{rmk}
%%   \todo{revise}
%%   Although the category of Tate objects has been considered by various authors, we believe that it is the category of Artin objects which is more tractible to understand as a first step towards understanding the Artin--Tate category. In particular, the first major step in understand both the $\C$ and $\R$ cases is the identification of the category of Artin objects via the functors $c$ and $c_{\C/\R}$.
%% \end{rmk}
%% \todo{There's a remark I commented here about how this says you should think about Artin before Tate.}

\begin{rmk}
  When working over a general field one might think to replace the Betti realization by some variant of etale localization (see \cite{EldenJay}). The analog of this theorem would be that the category of Artin objects is already etale local.
  However, in many examples (such as finite fields) one consequence of the Morel connectivity theorem \cite{MorelConn} is that this is not true. An important precursor to answering \Cref{qst:main} will be producing the variant of equivariant homotopy theory which appears as the category of Artin objects.
\end{rmk}
  
\begin{rmk}
  Using the functors $c$ and $c_{\C/\R}$ the mapping spaces in $\SHC$ and $\SHR$ can be upgraded into mapping spectra and mapping $C_2$-spectra respectively. % See [notatins] for details.
\end{rmk}

%% Summarizing the contents of this section, we have diagrams

%% \begin{center}
%%   \begin{tikzcd}
%%     \SHC \ar[d, "\b"] & \SHR \ar[r, "{(-)[a^{-1}]}"] \ar[l, "(-)_\C"] \ar[d,"\b"] &  \Mod ( \SHR ; \Ss[a^{-1}] ) \ar[d,"\b"] & \SHC \ar[r, "\mathrm{Nm}_{\C/\R}"] \ar[d, "\b"] & \SHR \ar[d, "\b"] \\
%%     \Sp & \Sp_{C_2} \ar[l,"\Phi^e"] \ar[r, "\Phi^{C_2}"] & \Sp & \Sp \ar[r, "\mathrm{Nm}_e^{C_2}"] & \Sp_{C_2}
%%   \end{tikzcd}
%%   \begin{tikzcd}    
%%     \Sp_{C_2} \ar[rr, "\mathrm{id}_{\Sp_{C_2}}", bend left=20] \ar[r, "c_{\C/\R}"'] \ar[d,"\Phi^e"] & \SHR \ar[r, "\b"'] \ar[d,"(-)_\C"] & \Sp_{C_2} \ar[d,"\Phi^e"] &  \Sp_{C_2} \ar[rr, "\mathrm{id}_{\Sp_{C_2}}", bend left=20] \ar[r, "\nu"'] \ar[d,"\Phi^e"] & \SHR \ar[r, "\b"'] \ar[d,"(-)_\C"] & \Sp_{C_2} \ar[d,"\Phi^e"] \\
%%     \Sp \ar[rr, "\mathrm{id}_{\Sp}"', bend right=20] \ar[r, "c"] & \SHC \ar[r, "\b"] & \Sp & \Sp \ar[rr, "\mathrm{id}_{\Sp}"', bend right=20] \ar[r, "\nu"] & \SHC \ar[r, "\b"] & \Sp
%%   \end{tikzcd}
%% \end{center}

\begin{summary}\label{prop:comp}
  There are commuting diagrams
  \begin{center}
  \begin{tikzcd}
    \SHR \ar[r, "(-)_\C"] \ar[d,"\b"] &
    \SHC \ar[d, "\b"] &
    \Sp_{C_2} \ar[rr, "\mathrm{id}_{\Sp_{C_2}}", bend left=20] \ar[r, "c_{\C/\R}"'] \ar[d,"\Phi^e"] &
    \SHR \ar[r, "\b"'] \ar[d,"(-)_\C"] &
    \Sp_{C_2} \ar[d,"\Phi^e"] \\
    \Sp_{C_2} \ar[r,"\Phi^e"] &
    \Sp &
    \Sp \ar[rr, "\mathrm{id}_{\Sp}"', bend right=20] \ar[r, "c"] &
    \SHC \ar[r, "\b"] &
    \Sp ,
  \end{tikzcd}
  \end{center}
  
  and various functors considered in this section enjoy the following properties:
  \begin{enumerate}
  \item Each of $(-)_{\C}$, $\b$, $\Phi^e$, $c$ and $c_{\C/\R}$ is a symmetric monoidal left adjoint.
  \item Both $(-)_{\C}$ and $\Phi^e$ are right adjoints as well.
  %% \item Both $\nu$ functors are lax symmetric monoidal and commute with filtered colimits.
  %% \item Both norm functors are symmetric monoidal and commute with sifted colimits.
  \item All of the functors restrict to the categories of Artin--Tate objects and retain the properties listed in (1) and (2). In later sections we will almost exclusively deal with these restrictions so we will not use distinct notation for them.
  \item On Picard elements the functors in the digrams above behave as follows:
  \end{enumerate}  
  %% \nu(\Ss^{2n+1}) &\cong \o_\C^{2n+1,n}  \\        
  \begin{center}
    \begin{tabular}{llll}
      $ \b(\o_\R^{p,q,w}) \cong \o_{C_2}^{p+q\sigma} $ &
      $ \b(\o_\C^{s,w}) \cong \Ss^s$ &
      $ c_{\C/\R}(\o_{C_2}^{p+q\sigma}) \cong \o_\R^{p,q,0}$ \\
      $ c(\Ss^s) \cong \o_\C^{s,0}$  &
      $ (\o_\R^{p,q,w})_\C \cong \o_\C^{p+q,w}$ &
      $ \Phi^e(\o_{C_2}^{p+q\sigma}) \cong \Ss^{p+q}$ &
      %$ \Phi^{C_2}(\o_{C_2}^{p+q\sigma}) \cong \Ss^{p}$ 
    \end{tabular}                  
  \end{center}
\end{summary}


  %% \begin{center}
  %%   \begin{tabular}{llll}
  %%     $ \b(\o_\R^{p,q,w}) \cong \o_{C_2}^{p+q\sigma} $ &
  %%     $ \b(\o_\C^{s,w}) \cong \Ss^s$ &
  %%     $ c_{\C/\R}(\o_{C_2}^{p+q\sigma}) \cong \o_\R^{p,q,0}$ &
  %%     $ c(\Ss^s) \cong \o_\C^{s,0}$ \\      
  %%     $ (\o_\R^{p,q,w})_\C \cong \o_\C^{p+q,w}$ &
  %%     $ \Phi^e(\o_{C_2}^{p+q\sigma}) \cong \Ss^{p+q}$ &
  %%     $ \o_\R^{p,q,w}[a^{-1}] \cong \o_\R^{p,0,w}[a^{-1}]$ &         
  %%     $ \Phi^{C_2}(\o_{C_2}^{p+q\sigma}) \cong \Ss^{p}$ \\
  %%     $ \nu(\o_{C_2}^{n\rho}) \cong \o_\R^{n,n,n}$ &
  %%     $ \nu(\Ss^{2n}) \cong \o_\C^{2n,n} $ &
  %%     $ \mathrm{Nm}_\C^\R(\o_\C^{s,w}) \cong \o_\R^{s,s,2w}  $ &
  %%     $ \mathrm{Nm}_e^{C_2}(\Ss^{n}) \cong \o_{C_2}^{n,n}$ 
  %%   \end{tabular}                  
  %% \end{center}


\begin{proof}
  %% We have already discussed all aspects of (1) except the functor $(-)[a^{-1}]$ whose behavior falls under \cite[Proposition 4.3.17]{ECII}.
  The only one of these which we have not already discussed is (2).
  The functor $(-)_\C$ can be described as tensoring up to $Ca$.
	Since $Ca$ is dualizable this functor commutes with all limits and colimits. Similarly, we may conclude that $\Phi^e$ commutes with all limits and colimits.
  %% Part (3) and the other statements about the functor $\nu_\R$ will be proved in \Cref{sec:nu}.
  %% Part (4) is a restatement of \cite[Proposition 4.5]{norms}.
  %% The only part of (5) which is not immediate is the claim about norm functors.
  %% This follow from the fact that $\mathrm{Nm}_{\C}^\R$ sends the bigraded spehres to trigraded spheres and the cross effect functor takes Artin--Tate objects to Artin--Tate objects. Finally, (6) is a collation of identifications we have already justified.
\end{proof}


\subsection{$\R$-motivic spectra as a deformation}\ 

In this subsection we refine the statement of \Cref{thm:main} into a sequence of precise claims which we will verify across the remainder of the paper. We summarized both \Cref{thm:C-main} and \Cref{thm:main} by saying that the category of Artin--Tate motivic spectra is a 1-parameter deformation of a purely topological category. To start, we clarify this, first over $\C$ and then over $\R$.

\begin{enumerate}
\item There is a distinguished element of the $p$-complete homotopy groups over $\C$,
  $\tau \in \pi_{0,-1}\Ss_p$, which maps to $1$ under Betti realization \footnote{The Betti realization map $\pi^{\C}_{0,-1}\Ss_p \to \pi_0\Ss_p$ is an isomorphism so the latter property uniquely identifies $\tau$.}. \cite[Lemma 23]{HKO}
\item The Betti realization functor factors through the category of $\tau$-local objects and provides a symmetric monoidal equivalence
  \[ \Mod ( \SHCat ; \Ss_p[\tau^{-1}] ) \simeq \Sp_{ip}. \]
  The idea for this goes back to \cite[Section 2.6]{DI}, but was first proven in \cite{Pstragowski}.  
\item The category $\SHCat$ can be equipped with the structure of a $\Sp_{ip}^{\Fil}$-algebra.
  More explicitly, this means we have a symmetric monoidal left adjoint
  \[ i^* : \Sp_{ip}^{\Fil} \to \SHCat, \]
  which sends the shift map in $\Sp_{ip}^{\Fil}$ to $\tau$. As a consequence of this $C\tau$ acquires the structure of a commutative algebra, a fact originally proven by Gheorghe \cite{Ctau}.
\item The category of modules over $C\tau$ is equivalent to a renormalization of the derived category of even $\BP_*\BP$-comodules. More specifically we have
  \[ \Mod ( \SHCat ; C\tau ) \simeq \IndCoh ( \Mfg )_{ip}. \]
  This equivalence sends $C\ta \otimes \Ss^{s,w}$ to $\Sigma^{s-2w} \omega_{\mathbb{G}/\Mfg} ^{\otimes w}$, 
  and on the level of homotopy groups it induces an equivalence
    \[ \pi_{s,w}^\C(C\tau) \cong \Ext_{\BP_*\BP}^{2w-s,w}(\BP_*, \BP_*). \] \todo{There is a small $p$-completion issue here...}
    The equivalence of categories is due to Gheorghe--Wang--Xu \cite{ctauactalol}, though the above isomorphism of groups was first proven by Isaksen \cite[Proposition 6.2.5]{StableStems} and then upgraded to an isomorphism of rings with all higher structure by Gheorghe \cite{Ctau}.
  \item There is an equivalence of symmetric monoidal categories between $\SHCat$ and $\Syn_{\MU, ip}^{\even}$, where the latter is the category of \emph{even} $\MU$-\emph{synthetic spectra} constructed in \cite{Pstragowski}. Notably, Pstr\k{a}gowski's construction uses no algebraic geometry, and requires only knowledge of the homotopy commutative ring object $\MU$ in $\Sp$. The comparison between these two categories proceeds via the close relationship between $\MGL$ and $\MU$. %first worked out by Levine. \todo{Was it?}
\item The adjunction $i$ is affine in the sense that there is a commutative algebra $R^{\C}_\bullet$ in $\Sp_{ip}^{\Fil}$ and an equivalence of symmetric monoidal categories under $\Sp_{ip}^{\Fil}$
  \[  \SHCat \simeq \Mod ( \Sp_{ip}^{\Fil} ; R^{\C}_\bullet ). \]
  Moreover, the commutative algebra $R^\C_\bullet$ admits an explicit construction which uses no algebraic geometry. It is given by
  \[ R_\bullet^{\C} \coloneqq \mathrm{Tot}^* \left( \tau_{\geq 2\bullet} \MU_p^{\otimes * + 1} \right). \]
    This construction uses only the commutative algebra $\MU$ in $\Sp$, and again the comparison proceeds via the close relationship between $\MGL$ and $\MU$. This approach is due to Gheorghe--Isaksen--Krause--Ricka \cite{cmmf}.
\end{enumerate}

%The key to understanding why we view $\SHCat$ as a 1-parameter deformation is point (2).
%Geometrically, a $1$-parameter deformation is a family over $\AA^1$.
%At the level of categories this means having a $\QCoh(\AA^1)$-algebra structure.
%Certain $1$-parameter deformations come with the extra structure of compatible isomorphisms relating the fiber over $\lambda$ and the fiber over $a\lambda$ for $a$ invertible.
%In the language of stacks we may describe this as having a family over $\AA^1/\G_m$.
%At the level of categories this means having a $\QCoh(\AA^1/\G_m)$-algebra structure.
%Finally, we recall that in spectral algebraic geometry there is an equivalence,
%$\QCoh(\AA^1/\G_m) \simeq \Sp^{\Fil}$ \footnote{By base-change a similar results holds over any base.}.
%Tracking through the various maps one further finds that the coordinate on $\AA^1$ corresponds to the shift map on filtered objects.

%In practice deformations typically come in two flavors.
%The first is where we examine how a central fiber of interest can deform (often over a formal base).%
%Visually one might image a central object spreading outwards.
%An example of this is the picture suggested by the Bogomolov--Tian--Todorov theorem on the unobstructedness of deformations of Calabi-Yau varieties over $\C$.
%The second is where we see a family of objects of interest degenerating inward to some special fiber. An example of this is the picture suggested by an elliptic fibration, with a family of elliptic curves degenerating towards a special point where we get a singular curve.
%Our situation of interest is firmly of the second type.
%Thus, if we wish to be more specific $\SHCat$ provides a degeneration of $\Sp_{ip}$ inward to an algebraic special fiber. More than just algebraic, the special fiber is algorithmic, i.e. any specific question about finite objects can be answered by a computer in finite time.

In fact, every aspect of the picture over $\C$ extends to $\R$ in the simplest reasonable way.
The reader who has previously studied $\R$-motivic spectra may find this rather surprising (as the authors did), and we suggest that this highlights the primacy of the Artin--Tate category over the Tate category.

\begin{enumerate}
\item There is a distinguished element of the $p$-complete homotopy groups over $\R$,
  $\ta \in \pi_{0,0,-1}\Ss_p$, which maps to $1$ under Betti realization \footnote{As above, the Betti realization map $\pi^{\R}_{0,0,-1}\Ss_p \to \pi^{C_2}_0\Ss_p$ is an isomorphism so $\ta$ is uniquely identified.}. We will construct $\ta$ in \Cref{sec:ta}.
\item The Betti realization functor factors through the category of $\ta$-local objects, and provides a symmetric monoidal equivalence
  \[ \Mod ( \SHRat ; \Ss_2[\ta^{-1}] ) \simeq \Sp_{C_2,i2}. \]
  This will be proven as the main theorem of \Cref{sec:homology}.  
\item The category $\SHRat$ can be equipped with the structure of a $\Sp_{i2}^{\Fil}$-algebra.
  More explicitly, this means we have a symmetric monoidal left adjoint
  \[ i^* : \Sp_{i2}^{\Fil} \to \SHRat \]
  which sends the shift map in $\Sp_{i2}^{\Fil}$ to $\ta$.
  As a consequence of this $C\ta$ acquires the structure of a commutative algebra.
  We may also regard $\SHRat$ as a $\Sp_{C_2}$-algebra via the functor $c_{\C/\R}$.
  Tensoring these two functors together we obtain a symmetric monoidal left adjoint
  \[ i^* : \Sp_{C_2, i2}^{\Fil} \to \SHRat. \]
    These statements will be proven in \Cref{sec:top} as part of \Cref{prop:top-diagram}.
\item The category of modules over $C\ta$ is equivalent to a renormalization of the derived category of an abelian category of equivariant $\BP_*\BP$-comodules. More precisely, we have an equivalence of presentably symmetric monoidal categories
  \[ \Mod ( \SHRat ; C\ta ) \simeq \Mod(\Sp_{C_2} ; \uZt) \otimes_{\Z} \IndCoh ( \Mfg ). \]
  This equivalence sends $C\ta \otimes \Ss^{p,q,w}$ to $\Sigma^{(p-w)+(q-w)\sigma} \uZt \otimes \omega_{\mathbb{G}/\Mfg} ^{\otimes w}$ and on the level of tri-graded rings of homotopy groups it induces an isomorphism \footnote{If one wants the tensor product can be moved outside the Ext, but then it must be taken in a derived sense.}
  \[ \pi_{p,q,w}^\R C\ta \cong \bigoplus_{w+a-s = p} \Ext_{\BP_*\BP}^{s,2w}(\BP_*, \BP_* \otimes \pi_{a + (q-w) \sigma}^{C_2} \uZt). \]
  This will be proven as the main theorem of \Cref{sec:cta}.
\item The adjunction $i$ is affine in the sense that there is a commutative algebra
  $R^{\R}_\bullet$ in $\Sp_{C_2,i2}^{\Fil}$ and an equivalence of presentably symmetric monoidal categories under $\Sp_{C_2,i2}^{\Fil}$
  \[ \SHRat \simeq \Mod ( \Sp_{C_2,i2}^{\Fil} ; R^{\R}_\bullet ). \]
  Moreover, the commutative algebra $R^\R_\bullet$ admits an explicit construction which uses no algebraic geometry. It is given by
  \[ R_\bullet^{\R} \coloneqq \mathrm{Tot}^* \left( P_{2\bullet} \MU_{\R,2}^{\otimes * + 1} \right), \]
  where $P_{n}$ is the functor which takes the $n^{\mathrm{th}}$ slice cover of a $C_2$ spectrum.
    This construction uses only the commutative algebra $\MU_\R$ in $\Sp_{C_2}$ introduced by Landweber \cite{LandweberI, LandweberII}. %\todo{This seems to be the original cite... or did you want something more modern?}
  The comparison proceeds via understanding the close relationship between $\MGL$ and $\MU_\R$ over $\R$.
  This will be proven as the main theorem of \Cref{sec:top}.
\end{enumerate}


%% \begin{prop}
%%   There is an equivalence of categories,
%%   \[ \SHR[2^{-1}] \simeq \Sp[2^{-1}] \times \SHC[2^{-1}]. \]
%% \end{prop}

%% \begin{proof}
%%   An idempotent in $\pi_0\o$ produces a product decomposition of the category.
%%   After inverting 2 we have $\left( \frac{[C_2]}{2} \right)^2 = \frac{[C_2]}{2}$.
%%   This splits the category $\SHR[2^{-1}]$ into two pieces.
%%   In one summand we have $[C_2]=0$ and $\rho$ is invertible.
%%   In the other summand $[C_2]$ is invertible.
%%   In \cite{Tom?}, Bachmann shows that $\SHR[2^{-1}, \rho^{-1}] \simeq \Sp$.

%%   We identify the category where $2$ and $[C_2]$ are invertible as follows.
%%   Using the fact that $c_{\C/\R}$ is fully faithful we can use some standard equivariant tricks.
%%   The object $\Spec(\C)$ which corresponds to $C_{2+}$ is self-dual and if let $p : \Spec(\C) \to \Spec(\R)$ be the standard projection map, then the map $[C_2]$ can be factored as the composition
%%   \[ \Ss \xrightarrow{\mathbb{D}p} \Spec(\C) \xrightarrow{p} \Ss. \]
%%   From equivaraint homotopy theory we know that $\mathbb{D}p \circ p = 2$.
%%   Thus, we have an equivalence $\Ss[2^{-1}, [C_2]^{-1}] \simeq \Spec(\C)[2^{-1}]$.
%%   Note that $[C_2]=2$ as an endomorphism of $\Spec(\C)$ so inverting $[C_2]$ has no further effect once we've inverted $2$.
%%   Using the equivalence between $\Mod ( \SHR ; \Spec(\C))$ and $\SHC$ discussed in the previous subsection we obtain the desired indentification.
%% \end{proof}



Using the second special element of the tri-graded homotopy groups of the sphere, $a \in \pi_{0,-1,0}^\R \Ss$, we can delve further into the structure of $\SHRat$. In order to do this we begin with a digression on the element $a_{\sigma}$ in $C_2$-equivariant homotopy theory.

\begin{prop}
  The category $\Sp_{C_2}$ can be viewed as a 1-parameter family with coordinate $a_\sigma$, special fiber $\Sp$ and generic fiber $\Sp$ in the sense that:
  \begin{itemize}
  \item The category of $a$-local objects can be identified with spectra, i.e.
    there is an equivalence of presentably symmetric monoidal categories
    \[ \Mod ( \Sp_{C_2} ; \Ss[ a_{\sigma}^{-1} ] ) \simeq \Sp. \]
  \item The cofiber of $a_\sigma$, which we will denote $Ca_\sigma$, can be endowed with a commutative algebra structure and there is an equivalence of presentably symmetric monoidal categories
    \[ \Mod ( \Sp_{C_2} ; Ca_{\sigma} ) \simeq \Sp. \]
    \item There is a monoidal left adjoint
      \[ i^* : \Sp^{\Gr} \to \Sp_{C_2}, \]
      which sends $\o(1)$ to $\Ss^{\sigma}$ \footnote{The authors will return to the question of whether this functor can be upgraded to a monoidal functor in from filtered spectra in a future work.}. Moreover, this adjunction is affine in the sense that we have an equivalence of categories
      \[ \Sp_{C_2} \simeq \Mod ( \Sp^{\Gr} ; R_\bullet^{C_2}), \]
      where $R_n^{C_2} \simeq \Sigma\R \mathrm{P}^{-n-1}_{-\infty}$.
  \end{itemize}  
\end{prop}

This result is certainly well known to experts. %\todo{Cite who?} \todo{X to doubt.}
We give a proof in \Cref{app:boilerplate} as Examples \ref{exm:c2-underlying} and \ref{exm:c2-graded}.
Since therein the pair of functors $\Phi^e$ and $\Phi^{C_2}$ out of $\Sp_{C_2}$ become identified with modding out by $a_\sigma$ and inverting $a_\sigma$ respectively, this proposition produces a diagram of symmetric monoidal left adjoints
\begin{center}
  \begin{tikzcd}
    \Sp & & \Sp_{C_2} \ar[ll, "\Phi^{C_2}"']  \ar[dl, "{(-)[a_\sigma^{-1}]}"] \ar[dr, "Ca \otimes -"'] \ar[rr, "\Phi^e"] & & \Sp \\
     & \Mod ( \Sp_{C_2} ; \Ss[a_\sigma^{-1}] ) \ar[ul, "\simeq"'] & & \Mod ( \Sp_{C_2} ; Ca ) \ar[ur, "\simeq"] . & 
  \end{tikzcd}
\end{center}

We are now free to use the functor $c_{\C/\R}$ to push our description of $\Sp_{C_2}$ as a 1-parameter deformation into $\SHRat$ and obtain a description of that category as a 2-parameter deformation of $\Sp_{i2}$ with coordinates $\ta$ and $a$. For example, we can give $Ca$ the structure of a commutative algebra since $c_{\C/\R}(a_\sigma) = a$. Note that by \Cref{prop:comp} the commutative algebra structure obtained in this way is the same as the one coming from the equivalence $Ca \simeq \Spec(\C)$. 

While the 1-parameter deformations considered up to now had only a special and a generic fiber, when considered as a 2-parameter deformation $\SHRat$ has various ``limiting behavoirs'' which we presently study. Since each parameter can be set to either $0$ or $1$,
or left unspecified, there are 9 total categories of interest.
We summarize what we know in the following table:

{\renewcommand{\arraystretch}{1.4}
\begin{center}
  \begin{tabular}{|c|c|c|c|c|} \hline
    & unspecified & $\ta = 0$ & $\ta = 1$ \\\hline
    unspecified & $\SHRat$ & $\Mod(\Sp_{C_2, i2} ; \uZt) \otimes_{\Z} \IndCoh ( \Mfg )$ & $\Sp_{C_2,i2}$ \\\hline
    $a = 0$ & $\SHC^{\at}_{i2}$ & $\IndCoh ( \Mfg )_{i2}$ & $\Sp_{i2}$ \\\hline
    $a = 1$ & ? & $\Mod(\Sp ; \F_2[u_{2\sigma}]) \otimes_\Z \IndCoh ( \Mfg )$ & $\Sp_{i2}$ \\\hline
  \end{tabular}
\end{center}
}

%% \begin{center}
%%   \begin{tabular}{|c|c|c|c|c|} \hline
%%     & unspecified & $\ta = 0$ & $\ta = 1$ \\\hline
%%     un & $\SHRat$ & $\uZ_p$-module $\BP_*\BP$-comodules & $\Sp_{C_2}$ \\\hline
%%     $a = 0$ & $\SHCat$ & $\BP_*\BP$-comodules & $\Sp$ \\\hline
%%     $a = 1$ & ? & $\F_2[u^2]$-module $\BP_*\BP$-comodules & $\Sp$ \\\hline
%%   \end{tabular}
%% \end{center}

The only identification in this table which we have not yet discussed is the one in the bottom middle.
However, one can easily obtain this from the identification of $C\ta$-modules and \Cref{lem:tensor-mod}:
\begin{align*}
  \Mod( \SHRat ; C\ta[a^{-1}] ) &\simeq \Sp_{i2} \otimes_{\Sp_{C_2,i2}} \Mod( \SHRat ; C\ta ) \\
  &\simeq \Sp_{i2} \otimes_{\Sp_{C_2,i2}} \Mod(\Sp_{C_2,i2} ; \uZt) \otimes_{\Z} \IndCoh ( \Mfg ) \\
  &\simeq \Mod(\Sp_{i2} ; \Phi^{C_2} \uZt) \otimes_{\Z} \IndCoh ( \Mfg ) \\
  &\simeq \Mod(\Sp ; \F_2[u_{2\sigma}]) \otimes_{\Z} \IndCoh ( \Mfg ) ,
\end{align*}
where the final step is just the identification of $\Phi^{C_2}\uZ$ with $\F_2[u_{2\sigma}]$ where $u_{2\sigma}$ is a polynomial generator in degree 2. As a corollary, we can give a description of the tri-graded homotopy groups of several objects in $\SHRat$ in terms of better understood categories.

\begin{cor}
  There are isomorphisms of rings:
  \begin{enumerate}
  \item $\pi_{p,q,w}^\R(Ca) \cong \pi_{p+q,w}^\C\Ss $,
  \item $\pi_{p,q,w}^\R(Ca \otimes C\ta) \cong \Ext_{\BP_*\BP}^{2w-p-q,w}(\BP_*, \BP_*) $,
  \item $\pi_{p,q,w}^\R(Ca[\ta^{-1}]_p) \cong \pi_{p+q}\Ss_p $,
  \item $\pi_{p,q,w}^\R(C\ta) \cong \bigoplus_{w+a-s = p} \Ext_{\BP_*\BP}^{s,2w}(\BP_*, \BP_* \otimes \pi_{a + (q-w) \sigma}^{C_2} \uZt)$,
  \item $\pi_{p,q,w}^\R(C\ta[a^{-1}]) \cong \bigoplus_{w+2a-s = p} \left( \F_2\{u^{2a}\} \otimes_{\F_2} \Ext_{\BP_*\BP}^{s,2w}(\BP_*, \BP_*/2)\right) $,    
  \item $\pi_{p,q,w}^\R(\Ss_2[\ta^{-1}]) \cong \pi_{p,q}^{C_2}\Ss_2$,
  \item $\pi_{p,q,w}(\Ss_2[a^{-1}, \ta^{-1}]) \cong \pi_{p}\Ss_2$.
  \end{enumerate}  
\end{cor}

We now turn to the category of $a$-local objects, only labelled ``?'' in the table above. It is one of the most mysterious actors in the story told in this paper and we shall devote \Cref{sec:a-local} to its study. As we shall see, it records a deformation of the category of spectra sitting between that which corresponds to the Adams--Novikov spectral sequence and that which corresponds to the classical Adams spectral sequence. The role it plays in mediating between these two spectral sequences remains to be understood.

Since $a = c_{\C/\R}(a_\sigma)$ and $a_\sigma = \b (a) $, we can construct a diagram of symmetric monoidal left adjoints
\begin{center} \begin{tikzcd}
    \SHR \ar[r, "{(-)[a^{-1}]}"] \ar[d,"\b"] &
    \Mod ( \SHR ; \Ss[a^{-1}] ) \ar[d,"\b"] \\    
    \Sp_{C_2} \ar[r, "\Phi^{C_2}"] &
    \Sp .
\end{tikzcd} \end{center}
%% where the only functors we have yet to discuss are $a$-localization and geometric fixed points.
%% Since geometric fixed points for $C_2$ can be described as localization at $a_\sigma$ and $\b(a) = a_\sigma$ we obtain the desired square of functors.
%% As far as the authors are aware this is an entirely new catgory which  We will devote \Cref{sec:a-local} to gaining a rough understanding of this category. 
%% In $C_2$-equivariant homotopy theory, inverting the class $a_\sigma \in \pi_{-\sigma} ^{C_2} \Ss$ corresponds to applying the the geometric fixed points functor $\Phi^{C_2} : \Sp_{C_2} \to \Sp$.
As the diagram shows, the category $\Mod(\SHR; \Ss_2 [a^{-1}])$ can be viewed as the target category of an ``$\R$-motivic geometric fixed points'' functor.
Since $\Phi^{C_2} (\MU_\R) \simeq \MO$, one might guess that $\Mod(\SHRat; \Ss_2 [a^{-1}])$ is related to Pstragowski's synthetic category $\Syn_{\MO}$ \footnote{The spectrum $\MO$ is a sum of copies of $\F_2$, and as we shall see this implies $\Syn_{\MO} \simeq \Syn_{\F_2}$.}. Indeed, we construct a comparison functor.

\begin{prop}
  There is a symmetric monoidal left adjoint
  \[\Re_{\F_2} : \Mod(\SHRat; \Ss_2 [a^{-1}]) \to \Syn_{\F_2}\]
  which sends $\Ss^{p,q,w}$ to $\Ss^{p,w}$.
\end{prop}

%% This category may be viewed as a $1$-parameter deformation of $\Sp_{i2}$ with parameter $\ta$. 
On the other side, an interesting functor into $\Mod(\SHRat; \Ss_2 [a^{-1}])$ may be constructed using the Bachmann--Hoyois motivic norm functors.

\begin{rec} \label{rec:norms}
  In \cite{norms}, Bachmann and Hoyois construct symmetric monoidal norm functors along finite \'etale maps.
  These norm functors may be thought of as an indexed tensor product and in the case of $\R$ they fit into the following diagram showing compatability with the Hill--Hopkins--Ravenel norm functors in equivariant homotopy theory \cite[Section 11]{norms}:
  \begin{center}
  \begin{tikzcd}
  \SHC \ar[r, "\mathrm{Nm}_{\C}^\R"] \ar[d, "\b"] & \SHR \ar[d, "\b"] \\
  \Sp \ar[r, "\mathrm{Nm}_e^{C_2}"] & \Sp_{C_2}.
  \end{tikzcd}
  \end{center}
  From \cite[{Example 3.5, Lemma 4.4, Example 4.10}]{norms} we can conclude that on bigraded spheres $\mathrm{Nm}_\C^{\R}(\Ss_\C^{s,w}) \simeq \Ss^{s,s,2w}$. The functor $\mathrm{Nm}_{\C}^\R$ is polynomial of degree 2 and when applied to a direct sum we have a formula \cite[Corollary 5.13]{norms}
  \[ \mathrm{Nm}_{\C}^\R(X \oplus Y) \simeq \mathrm{Nm}_{\C}^\R(X) \oplus i_*(X \otimes Y) \oplus \mathrm{Nm}_{\C}^\R(Y). \]
  Since $\mathrm{Nm}_{\C}^\R$ commutes with sifted colimits, it follows from the above that it restricts to a functor on the categories of Artin--Tate objects
  \[\mathrm{Nm}_{\C}^\R : \SHR^{\at} \to \SHC^{\at}.\]
\end{rec}

While the norm functor $\mathrm{Nm}_{\C}^\R$ is not exact, it becomes exact after $a$ is inverted.

\begin{prop}
  The composite
  \[ \SHC_{i2}^{\at} \xrightarrow{\mathrm{Nm}_\C^\R} \SHRat \xrightarrow{(-)[a^{-1}]} \Mod ( \SHRat ; \Ss_2[a^{-1}] ) \]
  is a symmetric monoidal left adjoint.
  On Picard elements this composite sends $\Ss_2^{s,w}$ to $\Ss_2^{s,*,2w}[a^{-1}]$.
  The map $\tau$ is sent to $\ta^2$.
\end{prop}

In conclusion, we have a pair of functors 
\[\SHC^{\at}_{i2} \xrightarrow{\mathrm{Nm}_\C^\R(-)[a^{-1}]} \Mod( \SHRat; \Ss_2 [a^{-1}]) \xrightarrow{\Re_{\F_2}} \Syn_{\F_2},\]
which are each the identity on spectra after inverting $\ta$ and $\tau$.
%% After inverting $\ta$ and $\tau$, both functors are identified with the identity on $\Sp_{i2}$.
%% In conclusion, we see that the category $\Mod(\SHRat; \Ss_2 [a^{-1}])$
%% encodes information intermediate between the classical Adams and Adams-Novikov spectral sequences.
It would be very interesting to obtain a better computational understanding of the $a$-local category and the behavior of these two functors. More ambitiously, we ask:
%% $\Mod(\SHRat; \Ss_2 [a^{-1}])$.

\begin{qst}
  Is there a good description of the full subcategory of $a$-local objects in $\SHRat$?
  What geometric information does the $a$-localization of a smooth projective variety $X$ remember?
\end{qst}

%% AAAAAAAAAAAAAAAAAAAAAAAAAAAAAAAAAAAAAAAAAAAAAAAAA

The careful reader will have noticed that we switched from $p$-completing when discussing $\C$ to $2$-completing when discussing $\R$. The reason for this is that at an odd prime the category $\SHR^{\at}_{ip}$ admits the following simple description in terms of $\SHCat$.
%% the following simple description of the category of $\R$ motivic spectra after inverting 2.

\begin{prop}
  There is an equivalence of categories
  \[ \SHR^{\at}_{ip} \simeq \Sp_{ip} \times \SHCat \times \SHCat. \]
\end{prop}

This will be proven in \Cref{sec:odd}.
%% We now return to studying the $i2$-complete category and will make no further mention of odd primes outside of  where we prove this proposition.




%% It is an important that the $p$-complete cellular $\C$-motivic category can be viewed as a deformation of $p$-complete spectra along the $p$-complete complex cobordism spectrum $\MU$ \cite{Pstragowski, cmmf, ctauactalol}. One of the key inputs to this result is the identification of a class $\tau \in \pi_{0,-1} ^{\C} \Ss_p$ with certain key properties recalled below.
%% \todo{This paragraph should probably appear elsewhere.}

%% %% \begin{rec} \label{rec:tau}
%% %%   For all primes $p$, there is a class $\tau \in \pi_{0,-1} ^\C \Ss_p$ with the following properties:
%% %%   \begin{enumerate}
%% %%     \item Under Betti realization, $\tau$ goes to $1 \in \pi_{0} \Ss_p$.
%% %%     \item The induced functor $\SHC^{cell} _{ip} [\tau^{-1}] \to \Sp_{ip}$ is an equivalence of categories.
%% %%     \item The cofiber of $\tau$, $C\tau$, admits a unique $\mathbb{E}_\infty$-structure and $\MGL$-homology induces an equivalence of categories $\Mod_{C\tau} \simeq \Comod_{\MU_* \MU}$.
%% %%   \end{enumerate}
%% %% \end{rec}
%% %% \todo{Who to cite for this?}

%% %% The remainder of this section will be dedicated to \Cref{thm:ta} below, which constructs the real motivic analogue of $\tau$ and proves the analogue of (1) above.
%% %% We will address item (2) in \Cref{thm:comparison-invert-ta} and item (3) in \Cref{thm:cta}.

%% %
%% %\begin{dfn}
%% %  Let $a \in \pi_{0,-1,0}^\R$ denote the image of the stabilization of the $C_2$-equivariant map $a_\sigma : S^0 \to S^{\sigma}$ under the functor $c_{\C/\R}$.
%% %\end{dfn}
%% %
%% %\begin{cor}
%% %  There is a relation $\ta a = \rho$ in $\pi_{0,-1,-1} ^\R$.
%% %\end{cor}
%% %
%% %\begin{proof}
%% %  We have $\pi_{0,-1,-1} ^\R = \Z/2\{\rho\}$ by CITE, so it suffices to show that $\ta a \neq 0$. But $\b (\ta a) = a_\sigma \neq 0$, so this holds.
%% %\end{proof}
%% %\todo{Move this section.}

%% %% Our final major result is an explicit homotopy theoretic description of $\SHR^{gc}$, divorced from any algebro-geometric origin.

%% %% \begin{thm} \label{thm:filtered-model}
%% %% There is a filtered, $C_2$-equivariant $\mathbb{E}_\infty$-algebra $R_\bullet$ such that the category of filtered modules over $R_\bullet$ is equivalent to $\SHR^{gc}$.  Explicitly, the $n$th filtered piece $R_n$ of this ring is calculated as a totalization of a cosimplicial object with $k$th layer $P_{2n} \left((\MUR)^{\otimes k+1}\right)$.  Here, $P_{2n}$ refers to the $2n$th slice cover in $C_2$-equivariant homotopy theory, and $\MUR$ to the $C_2$-spectrum of Real bordism.
%% %% \end{thm}


%% This section is the beginning of the second part of this paper.
%% In the first part our main concern was providing a closed-form description of the galois-cellular category. In particular, we saw $\SHgc$ as a 2-parameter deformation of the category of spectra with coordinates $\ta$ and $a$. The second part of this paper we will be devoted to closely studying the structure of this category.
%% We begin by studying the limiting behaviors of this 2-parameter deformation.
%% Next, we investigate several $t$-structures on this category and their associated spectral sequences.
%% Finally, we will clarify the implications this has on the structure of the category $\SHgc$.
%% Throughout this section we pay particular attention to trigraded homotopy groups, using them as a yardstick for our understanding of various aspects of Artin--Tate $\R$-motivic spectra.

%% At certain points the proofs will become involved and when this happens we will defer a more complete discussion to its own section. The remaining NUMBER sections of the body of this paper are each of this type. For this reason, after reading this section the remaining sections of this paper can be read independently of one another.


%% In the previous subsection we saw $\SHRat$ as a 2-parameter deformation of the category of spectra with coordinates $\ta$ and $a$. In this section we study the ``limiting behavoirs'' of this deformation.
%% Throughout this section we pay particular attention to trigraded homotopy groups, using them as a yardstick for our understanding.

%% Before proceeding we must make the sense in which we use ``deformation'' precise.
%% In [CITE] we showed that $\SHgc$ is a 1-parameter deformation of $\Sp_{C_2}$ in a strong sense (the parameter $\ta$ arises via an action of $\Sp^{\Fil}$ on $\SHgc$). In [???], the authors show that $\Sp_{C_2}$ is a 1-parameter deformation of $\Sp$ with parameter $a$, albeit in a weaker sense. Given this mis-match of strucutre we spell out what we actually use in the following lemma.

%% \begin{lem}
%%   $\SHgc$ is a 2-parameter deformation of $\Sp$ with parameters $\ta$ and $a$ in the following sense.
%%   \begin{itemize}
%%   \item There are elements $\ta \in \pi_{0,0,-1}^\R$ and $a \in \pi_{0,-1,0}^\R$.
%%   \item Both $C\ta$ and $Ca$ naturally acquire the structure of commutative algebras.
%%   \item The category of modules over $\o [\ta^{-1}, a^{-1}]$ is equivalent to the category of spectra.
%%   \end{itemize}  
%% \end{lem}

%% \begin{proof}
%%   The elements $\ta$ and $a$ where constructed in [location].
%%   The commutative algebra structure on $C\ta$ was provided earlier in [location].
%%   In [location] we noted that $Ca \cong \Spec(\C)$. Since $\Spec(\C)$ is self-dual the diagonal provides it with the structure of a commutive algebra.
%%   In [location] we proved that $\o [\ta^{-1}]$-modules are equivalent to $C_2$-spectra. Now, it suffices to note that inverting $a$ in $C_2$-spectra gives ordinary spectra.
%% \end{proof}

%% Using this lemma,
%% each parameter can be set to either $0$ or $1$,
%% %\footnote{The value $\epsilon$ is supposed to be infinitesimal and correspond to the formal neighborhood of $0$.}
%% which together with ``unspecified'' gives 9 total commutative algebras whose category of modules we would like to study. These constitute the ``limiting behaviors'' of this 2-parameter deformation.
%% We summarize what we know in the following, rather cryptic, table.

%% \begin{center}
%%   \begin{tabular}{|c|c|c|c|c|} \hline
%%     & unspecified & $\ta = 0$ & $\ta = 1$ \\\hline
%%     un & $\SHgc$ & $\uZt$-module $\BP_*\BP$-comodules & $\Sp_{C_2}$ \\\hline
%%     %% $a$-cplt & ?\footnote{From the second column one can see that $C\ta$ is not $a$-complete. This implies that $\Ss$ is also not $a$-complete which stands in contrast to the $C_2$-equivariant setting.} & $\Z_2[a]$-module $\BP_*\BP$-comodules & $Fun(BC_2, \Sp)$ \\\hline 
%%     $a = 0$ & $\SH(\C)$ & $\BP_*\BP$-comodules & $\Sp$ \\\hline
%%     $a = 1$ & ?\footnote{This category is totally novel as far as we can tell.} & $\F_2[u^2]$-module $\BP_*\BP$-comodules & $\Sp$ \\\hline
%%   \end{tabular}
%% \end{center}

%% Leveraging the work from the first part of the paper most of what we have to say will follow rather formally from the results of the first part of the paper. However, the analysis of the category of $\Ss[a^{-1}]$-modules is more difficult and we will defer the proofs in this subsection to their own separate section. 

%% %% will be sufficiently complicated that it will require its own seperate section. This part is the analysis of the category of $\Ss[a^{-1}]$-modules. In the end we find that this category has various similarities to the $\F_2$-synthetic category of Pstragowski \cite{Pstragowski}. Moreover, a suitable enrichment of the Betti realization functor in fact lands in this category of $\F_2$-synthetic spectra. However, as part of the computations carried out in [appendic C] we will find that this comparison map does not admit a section. 

%% %% The first is the analysis of category of modules over $C\ta$.
%% %% As in the $\C$-motivic case we find that this category is algebraic in the sense that it is the derived category of an explicit abelian category which we will call ``Mackey functor $\BP_*\BP$-comodules''.

%% %% The second 

%% \begin{prop}
%%   Seven of the limiting behaviors from the table above admit direct descriptions.
%%   \begin{enumerate}
%%   \item The category of modules over $Ca$ is equivalent to
%%     $\SH(\C)$.
%%     Under this equivalence $\ta$ maps to $\tau$ and $Ca \otimes \o^{p,q,w}$ is sent to $\o^{p+q,w}$.
%%   \item The category of modules over $Ca \otimes C\ta$ is equivalent to
%%     $\mathrm{IndCoh}(\mathcal{M}_{\mathrm{fg}})$.
%%     Under this equivalence $Ca \otimes C\ta \otimes \o^{p,q,w}$ is sent to a copy of $\BP_*$ placed in
%%     topological degree $p+q$ and homological degree $2w - p - q$.    
%%   \item The category of modules over $Ca[\ta^{-1}]$ is equivalent to
%%     $\Sp$.
%%     Under this equivalence $Ca[\ta^{-1}] \otimes \o^{p,q,w}$ is sent to $\Ss^{p+q}$.
%%   \item The category of modules over $C\ta$ is equivalent to
%%     \[ \Mod ( \Sp_{C_2, i2}; \uZt) \otimes_{\Z} \mathrm{IndCoh}(\mathcal{M}_{\mathrm{fg}}). \]
%%     Under this equivalence $C\ta \o^{p,q,w}$ is sent to $(\Sigma^{(p-w)+(q-w)\sigma}\uZt) \otimes \BP_*$ where the copy of $\BP_*$ is placed in the heart and has internal degree $2w$.
%%   \item The category of modules over $C\ta[a^{-1}]$ is equivalent to
%%     \[ \Mod ( Vect_{\F_2}; \F_2[x] ) \otimes_{\Z} \mathrm{IndCoh}(\mathcal{M}_{\mathrm{fg}}) \]
%%     where $|x| = 2$.
%%     Under this equivalence $C\ta[a^{-1}] \otimes \o^{p,q,w}$ is sent to $\F_2[x] \otimes \BP_*$ where the copy of $\BP_*$ is placed in topological degree $p$ and homological degree $2w-p$.
%%   \item The category of modules over $\Ss[\ta^{-1}]$ is equivalent to
%%     $\Sp_{C_2}$.
%%     Under this equivalence $\o^{p,q,w}[\ta^{-1}]$ is sent to $\Ss^{p+q\sigma}$.
%%   \item The category of modules over $\Ss[\ta^{-1}, a^{-1}]$ is equivalent to
%%     $\Sp$.
%%     Under this equivalence $\o^{p,q,w}[\ta^{-1}, a^{-1}]$ is sent to $\Ss^{p}$.
%%   \end{enumerate}
%%   In each of the bullet points above ``equivalence'' is meant in the category of presentably symmetric monoidal categories. 
%% \end{prop}

%% Before proceeding to the proof of this proposition we extract information about the homotopy groups of the unit in each of these categories as a corollary.

%% \begin{cor}
%%   There are natural isomorphisms of rings,
%%   \begin{enumerate}
%%   \item $\pi_{p,q,w}^\R(Ca) \cong \pi_{p+q,w}^\C $,
%%   \item $\pi_{p,q,w}^\R(Ca \otimes C\ta) \cong \Ext_{\BP_*\BP}^{2w-s,w}(\BP_*, \BP_*) $,
%%   \item $\pi_{p,q,w}^\R(Ca[\ta^{-1}]) \cong \pi_{p+q} $,
%%   \item $\pi_{p,q,w}^\R(C\ta) \cong \bigoplus_{w+a-s = p} \Ext_{\BP_*\BP}^{s,2w}(\BP_*, \BP_* \otimes \pi_{a + (q-w) \sigma}^{C_2} \uZt)$,
%%   \item $\pi_{p,q,w}^\R(C\ta[a^{-1}]) \cong \left( \F_2[x] \otimes_{\F_2} \Ext_{\BP_*\BP}^{?,?}(\BP_*, \BP_*/2)\right)_{a,b} $,    
%%   \item $\pi_{p,q,w}^\R(\Ss[\ta^{-1}]) \cong \pi_{p,q}^{C_2}$,
%%   \item $\pi_{p,q,w}(\Ss[a^{-1}, \ta^{-1}]) \cong \pi_{p}$.
%%   \end{enumerate}  
%% \end{cor}

%% \begin{proof}[Proof (of \Cref{??}).]\ 
  
%%   (1) Since the commutative algebra structure on $Ca$ comes from the equivalence $Ca \cong \Spec(\C)$ we obtain the following proposition as a corollary of [cite].
%%     Since $a$ is in the image of $c$ the first equivalence follows from the fact that $c(C_2+) = \Spec(\C)_+$. The equivalence of $\Spec(\C)$-mod with $\SH(\C)$ follows from ???
%%   \todo{ Add 0-affineness of SH(l) / SH(k) into the appendix. }

%%   (2) is a corollary of (1) and [ctau paper].

%%   (3) is a corollary of (1) and [invert tau cite].

%%   (4) is just a restatement of the main results of \Cref{sec:cta}.

%%   (5) is obtained by substituting the equivalence
%%   \[ \Mod ( \Sp_{C_2, i2}; \uZt[a^{-1}]) \cong \Mod ( \Sp_{i2} ; \F_2[x] ) \cong \Mod( Vect_{\F_2} ; \F_2[x] ) \]
%%   where $|x| = 2$ into \Cref{thm:Cta}.

%%   (6) is just a restatement of \Cref{sec3 thing}.

%%   (7) is a corollary of (?). 
%%   In [location] we showed that the category of $\Ss[\ta^{-1}]$-modules is equivalent as a presentably symmetric monoidal category to $\Sp_{C_2}$. In [location], it is shown that the category of $\Ss[a^{-1}]$-modules in $\Sp_{C_2}$ is equivalent as a presentably symmetric monoidal category to $\Sp$.
%% From this we can derive the following corollary.

  
%% \end{proof}


%% \begin{qst}
%%   What is a good description of the full subcategory of $a$-local objects $\SHRat$.
%% \end{qst}


%% %% \begin{thm}
%% %%   As trigraded commutative rings, there are isomorphisms\footnote{If one wants the tensor product can be moved outside the Ext, but then it must be taken in a derived sense.}
%% %%   \[ \pi_{p,q,w}^\R(C\ta) \cong \bigoplus_{w+a-s = p} \Ext_{\BP_*\BP}^{s,2w}(\BP_*, \BP_* \otimes \pi_{a + (q-w) \sigma}^{C_2} \uZt). \]
%% %% \end{thm}


%% \subsection{Limiting behavoirs of the deformation}\ 

%% \input{Intro-limit.tex}

\subsection{Potpourri}\

In this subsection, we collect a number of additional results about the category of Artin--Tate $\R$-motivic spectra that are worth highlighting.%, but which didn't fit into the previous sections.

%% First we discuss the trigraded $\R$-motivic homology of a point, as well as the trigraded $\R$-motivic dual Steenrod algebra.
%% We then discuss the functor $\nur : \Sp_{C_2, i2} \to \SHRat$ and show that it may be used to recover many interesting Artin-Tate $\R$-motivic spectra. For more information on this, we refer the reader to \Cref{sec:nu}.
%% After that, we discusss the category of $a$-local Artin--Tate $\R$-motivic spectra and its relation to the Pstragowski's category of $\F_2$-synthetic spectra. \Cref{sec:a-local} discusses this topic in more detail.
%% Finally, we include some speculations on the possibility of other deformations of equivariant categories.

\subsubsection{Trigraded $\R$-motivic homology and Steenrod algebra}
From a computational point of view, an important first step in studying the trigraded homotopy groups of $\R$-motivic spectra is the computation of the trigraded homology of a point and the trigraded dual Steenrod algebra. Leaning on known bigraded computations, we make these computations in \Cref{sec:homology}.
%
%is then interested in determining the trigraded homology of a point, the motivic Steenrod algebra, and the trigraded homotopy groups of $C\ta$.  Leaning on known bigraded computations, we determine all of these invariants:

\begin{prop}
As trigraded commutative rings, with $|\ta|=(0,0,-1)$, there are isomorphisms
$$\pi^{\R}_{p,q,w} \HZt \cong \left(\pi^{C_2}_{p+q\sigma}\uZt\right)[\ta],$$
$$\pi^{\R}_{p,q,w} \HFt \cong \left(\pi^{C_2}_{p+q\sigma}\uFt\right)[\ta].$$
Here, $\pi^{C_2}_{p+q\sigma}\uZt$ and $\pi^{C_2}_{p+q\sigma}\uFt$ are placed in degrees $(p,q,0)$
\end{prop}

\begin{rmk}\label{rmk:ta-explanation}
  There are well-known elements $\tau \in \pi_{1,-1-1} ^{\R} \HFt$ and $u_{\sigma} \in \pi_{1-\sigma} ^{C_2} \uFt$. Write $u \in \pi_{1,-1,0} ^{\R} \HFt$ for the corresponding element.
  Then there is a relation (see \Cref{cor:rels})
  \[\tau = \ta \cdot u,\]
  i.e.
  \[\mathrm{tau} = \mathrm{ta} \cdot u.\]
  This was the original motivation for our choice of the character $\ta$.
\end{rmk}

\begin{prop} 
The trigraded dual Steenrod algebra $\pi^{\R}_{*,*,*} \left(\HFt \otimes \HFt\right)$ is isomorphic to
$$\left(\HFt\right)_{***} [\tau_0,\tau_1,\cdots,\xi_1,\xi_2,\cdots]/(\tau_i^2 = \ta (a \tau_{i+1} + u \xi_{i+1} + a \tau_0 \xi_{i+1}))$$
Here, $|\tau_i|=(2^i,2^{i}-1,2^{i}-1)$ and $|\xi_{i}|=(2^i-1,2^i-1,2^i-1)$.
\end{prop}

\begin{rmk}
The reader may notice that the `negative cone' in the $C_2$-equivariant homology of a point appears in the above formulas.  This is a feature of the Artin--Tate category, as opposed to the Tate category.
\end{rmk}

We expect that a computer could be coaxed into computing the $\HFt$-Adams spectral sequence for the tri-graded homotopy of $C\ta$, and that the natural maps between this spectral sequence and the $\HFt$-Adams spectral sequence for the tri-graded homotopy of the sphere would provide many of the differentials in the latter. This technique would likely be the most efficient way of computing both tri-graded $\R$-motivic stems and (in the long-run) the bi-graded $C_2$-equivariant homotopy groups. However, as the quad-graded nature of this computation would begin to tax our ability to visualize data it is unlikely that such a computation could be accomplished without the aid of a well-developed software suite for manipulating spectral sequence data.

%
%\begin{rmk}
%While the trigraded homotopy groups of $\HZt$ and $\HFt$ are $\ta$-torsion free, the trigraded homotopy groups of $\MGL$ are not.  This a key difference between the situations over $\C$ and $\R$, related to the fact in $C_2$-equivariant homotopy theory that the slice spectral sequence for $\MUR$ is non-trivial.
%\end{rmk}
%

\subsubsection{The effective slice filtration}
In \Cref{sec:eff}, we study Voevodsky's effective slice filtration. 
The effective slice filtration associates to any object $E \in \SHk$ a tower
\[ \dots \to f_{2} E \to f_1 E \to f_0 E \to f_{-1} E \to f_{-2} E \to \dots \to E.\]
The main result of the section is the following proposition, proven as \Cref{thm:eff-fil} which identifies the Betti realization of this tower.

\begin{prop}
  The functor $i_* : \SHRat \to \Sp_{C_2, i2}^{\Fil}$ sending $E$ to the tower
  \[ \dots \to \Map_{\SHRat} (\Ss_2 ^{0,0,1}, E) \xrightarrow{\ta} \Map_{\SHRat} (\Ss_2 ^{0,0,0}, E) \xrightarrow{\ta} \Map_{\SHRat} (\Ss_2 ^{0,0,-1}, E) \to \dots \]
  is equivalent to the functor $\b \circ f_\bullet : \SHRat \to \Sp_{C_2, i2}^{\Fil}$ taking $E \in \SHRat$ to the Betti realization of its effective slice tower.
\end{prop}

\begin{rec}
  Voevodsky has defined the ``rigid homotopy groups'' \cite[Definition 5.1]{VoeOpenProbs} of an $\R$-motivic spectrum $X$ to be
  \[\pi_{p,q,w} ^{\R,rig} X \coloneqq \pi_{p,q,w} ^{\R} s_w X,\]
  where $s_w X$ is $n$th effective slice of $X$, i.e. the cofiber of $f_{n+1} X \to f_n X$.
\end{rec}

Using the notion of rigid homotopy groups, he observed the phenomenon of \emph{algebraic degeneration} in motivic homotopy theory.
Since the associated graded of this tower also computes the homotopy groups of $X \otimes C\ta$, it follows from the proposition that $\pi_{p,q,w} ^{\R,rig} X \cong \pi_{p,q,w} ^{\R} X \otimes C\ta$ when $X \in \SHRat$.
This shows that Voevodsky's ideas about algebraic degeneration line up with notions coming from the cofiber of $\ta$ (or $\tau$) philosophy.

The proposition also implies that the so-called $C_2$-effective spectral sequence (see \cite{Kong}) is interchangeable with the $\ta$-Bockstein spectral sequence.

%The rigid homotopy groups introduced in \cite{VoeOpenProbs} can be described as the homotopy groups of the associated graded of this tower.
%Since the associated graded of this tower also computes the homotopy groups of $C\ta$ we learn that these two natural objects are equivalent.

\subsubsection{The functor $\nur$}
%
%Such a concrete description of $\SHgc$ allows one to build $\R$-motivic spectra `by hand,' using purely the tools of $C_2$-equivariant homotopy theory.
%
%We also consider two exotic comparison functors: $\nu_\R$ and the norm functor (introduced in \cite{norms}).
%These functors produce more interesting examples of objects of $\SHRat$ at the cost of being non-exact.
In \Cref{sec:nu}, we will construct a lax symmetric monoidal functor
\[\nur : \Sp_{C_2, i2} \to \SHRat,\]
which is a section of the Betti realization functor and sends $\Ss_2^{n\rho}$ to $\Ss_2^{n,n,n}$.
This functor is defined similarly to Pstragowski's synthetic analog functor \cite[Definition 4.3]{Pstragowski}.
Unlike the section $c_{\C/\R}$ of Heller--Ormsby, the functor $\nur$ is not exact.
However, it is better adapted to the construction of interesting $\R$-motivic spectra, as the following result shows.

\begin{prop}
  There are equivalences of commutative algebras,
  \begin{center}
  \begin{tabular}{lll}
    $ \nur \uFt \simeq \HFt,  \qquad\qquad$ &
    $ \nur \uZt \simeq \HZt, \qquad\qquad$ &    
    $ \nur \MU_{\R,2} \simeq \MGL_2, $ \\
    $ \nur \ku_{\R,2} \simeq \kgl_2, $ &
    $ \nur \ko_{C_2,2} \simeq \kq_2.$
  \end{tabular}
  \end{center}
\end{prop}

On the one hand, this proposition suggests that each of the $2$-complete $\R$-motivic homology theories above is not particularly far from being purely topological.
On the other hand, it suggests that one may profitably define $\R$-motivic analogs of $2$-complete $C_2$-equivariant homology theories by applying $\nu_\R$.
For example, if we had a good definition of $2$-complete $C_2$-equivariant connective topological modular forms $\tmf_{C_2, 2}$, then we could define a spectrum of $\R$-motivic modular forms as $\nu_{\R} \tmf_{C_2, 2}$. Unfortunately, no such definition is currently known.
%
%\begin{rec}
%  In \cite{norms}, Bachmann and Hoyois construct symmetric monoidal norm functors along finite etale maps.
%  These norm functors may be thought of as an indexed tensor product and in the case of $\R$ they fit into the following diagram showing compatability with the Hill--Hopkins--Ravenel norm functors in equivariant homotopy theory \cite[Section 11]{norms}
%  \begin{center}
%  \begin{tikzcd}
%  \SHC \ar[r, "\mathrm{Nm}_{\C/\R}"] \ar[d, "\b"] & \SHR \ar[d, "\b"] \\
%  \Sp \ar[r, "\mathrm{Nm}_e^{C_2}"] & \Sp_{C_2}
%  \end{tikzcd}
%  \end{center}
%  From \cite[{Example 3.5, Lemma 4.4, Example 4.10}]{norms} we can conclude that on bigraded spheres $\mathrm{Nm}_\C^{\R}(\Ss_\C^{s,w}) \simeq \Ss^{s,s,2w}$. The functor $\mathrm{Nm}_{\C}^\R$ is polynomial of degree 2 and when applied to a direct sum we have a formula \cite[Corollary 5.13]{norms}
%  \[ \mathrm{Nm}_{\C}^\R(X \oplus Y) \simeq \mathrm{Nm}_{\C}^\R(X) \oplus i_*(X \otimes Y) \oplus \mathrm{Nm}_{\C}^\R(Y). \]
%\end{rec}
%
%
%
%\begin{thm}
%There is an equivalence of $\R$-motivic spectra
%$$\nu(\mathrm{ko}_{C_2}) \simeq kq.$$
%\end{thm}
%
%We expect that the $\nu$ functor will be helpful in producing additional $\R$-motivic spectra of higher chromatic height.  The $\nu$ functor may also provide a computational tool for studying the $C_2$-equivariant homotopy groups of any $C_2$-spectrum $X$, via the comparison $\nu X \to \nu X \otimes C\ta$.










%% AAAAAAAAAAAAAAAAAAAAAAAAAAAAAAAAAAAAAAAAAAAAAAAAAAA

%% material from comparisons

%% AAAAAAAAAAAAAAAAAAAAAAAAAAAAAAAAAAAAAAAAAAAAAAAAAAA

%% We also consider two exotic comparison functors: $\nu_\R$ and the norm functor (introduced in \cite{norms}).
%% These functors produce more interesting examples of objects of $\SHRat$ at the cost of being non-exact.
%% The functor $\nu_\R$ will be constructed in [location] where we will prove the following proposition:

%% \begin{prop}
%%   There is commutative diagram of lax symmetric monoidal functors,
%%   \begin{center}
%%       \begin{tikzcd}    
%%         \Sp_{C_2} \ar[rr, "\mathrm{id}_{\Sp_{C_2}}", bend left=20] \ar[r, "\nu_\R"'] \ar[d,"\Phi^e"] & \SHR \ar[r, "\b"'] \ar[d,"(-)_\C"] & \Sp_{C_2} \ar[d,"\Phi^e"] \\
%%         \Sp \ar[rr, "\mathrm{id}_{\Sp}"', bend right=20] \ar[r, "\nu_\C"] & \SHC \ar[r, "\b"] & \Sp.
%%       \end{tikzcd}
%%   \end{center}
%%   where the functor $\nu_\C$ is the functor constructed in [location].
%%   These functors have the property that
%%   \begin{center}
%%   \begin{tabular}{lll}
%%     $ \nu \Ss_2^{n \rho} \simeq \Ss_2^{n,n,n}  \qquad\qquad$ &
%%     $ \nu \uFt \simeq M\F_2  \qquad\qquad$ &
%%     $ \nu \uZt \simeq M\Z_2 $ \\
%%     $ \nu \MU_{\R,2} \simeq \MGL_2 $ &
%%     $ \nu \ku_{\R,2} \simeq \kgl_2 $ &
%%     $ \nu \ko_{C_2,2} \simeq \kq_2$
%%   \end{tabular}
%%   \end{center}
%% \end{prop}

%% On the one hand, this proposition suggests that each of these motivic homology theories above is not particularly far from being purely topological.
%% On the other hand, this proposition suggest that one may profitably define $\R$-motivic analogs of $C_2$-equivariant homology theories by applying $\nu_\R$.
%% As an example which is still out reach at the moment, if we had a good definition of $\tmf_{C_2}$, then one could define $\R$-motivic modular forms as $\nu_{\R} \tmf_{C_2}$. 

%% \begin{rec}
%%   In \cite{norms}, Bachmann and Hoyois construct symmetric monoidal norm functors along finite etale maps.
%%   These norm functors may be thought of as an indexed tensor product and in the case of $\R$ they fit into the following diagram showing compatability with the Hill--Hopkins--Ravenel norm functors in equivariant homotopy theory \cite[Section 11]{norms}
%%   \begin{center}
%%   \begin{tikzcd}
%%   \SHC \ar[r, "\mathrm{Nm}_{\C/\R}"] \ar[d, "\b"] & \SHR \ar[d, "\b"] \\
%%   \Sp \ar[r, "\mathrm{Nm}_e^{C_2}"] & \Sp_{C_2}
%%   \end{tikzcd}
%%   \end{center}
%%   From \cite[{Example 3.5, Lemma 4.4, Example 4.10}]{norms} we can conclude that on bigraded spheres $\mathrm{Nm}_\C^{\R}(\Ss_\C^{s,w}) \simeq \Ss^{s,s,2w}$. The functor $\mathrm{Nm}_{\C}^\R$ is polynomial of degree 2 and when applied to a direct sum we have a formula \cite[Corollary 5.13]{norms}
%%   \[ \mathrm{Nm}_{\C}^\R(X \oplus Y) \simeq \mathrm{Nm}_{\C}^\R(X) \oplus i_*(X \otimes Y) \oplus \mathrm{Nm}_{\C}^\R(Y). \]
%% \end{rec}






\subsubsection{Further deformations}

One of our motivations for this project was the hope that a deformation-theoretic description of Artin--Tate $\R$-motivic stable homotopy theory would suggest other profitable deformations of classical stable homotopy theories.

In \Cref{app:def-const}, we show how a $1$-parameter deformation of a stable presentably symmetric monoidal category $\CC$ may be associated to a choice of object $E \in \CC$ and filtration $\{\CC_{\geq k}\}$ of $\CC$.
The case of $\SHCat$ is associated to the case where $E = \MU_{p}$, $\CC = \Sp_{ip}$ and $\CC_{\geq k}$ consists of the $2k$-connective objects. 
The case of $\SHRat$, on the other hand, is associated to the case where $E = \MU_{\R, 2}$, $\CC = \Sp_{C_2, i2}$ and $\CC_{\geq k}$ consists of the regular slice $2k$-connective objects.
Note that while the completions are necessary for the link to motivic homotopy theory, for the purpose of studying classical homotopy theory one may just as well work with integral deformations.

This suggests several other deformations that may be profitable to study.
\begin{enumerate}
\item We can take the deformation of $C_{2^n}$-equivariant homotopy theory with respect to the norm $\mathrm{N}_{C_2} ^{C_{2^n}} \MU_\R$ and the even slice filtration.
\item At odd primes, one could try to construct the deformation associated to the hypothetical spectrum $\mathrm{BP}_{\mu_p}$ and the even regular slice filtration.
\item One could take the deformation of $\Sp_{C_p}$ associated to $\underline{\mathbb{F}}_p$ and the regular slice filtration.
At the prime $2$, this is connected to a variant of the $\underline{\mathbb{F}}_2$-Adams spectral sequence.
Dylan Wilson has suggested that this variant may be more tractable at odd primes than the $\underline{\mathbb{F}}_p$-Adams spectral sequence, considering his forthcoming work with Krishanu Sankar proving that $\underline{\mathbb{F}}_p \otimes \underline{\mathbb{F}}_p$ is not a free $\underline{\mathbb{F}}_p$-module.
%Although the $C_p$-equivaraint Steenrod algebra is not yet totally known, the deformation of $\Sp_{C_p}$ associated to $\underline{\mathbb{F}}_p$ and the regular slice filtration, may be tractable anyway.
%This may be particularly interesting when $p$ is odd, as in this case the regular slice tower for $\underline{\mathbb{F}}_p \otimes \underline{\mathbb{F}}_p$ is nontrivial.
\item The deformation of $\Syn_{\MU}$ based on $\nu \F_p$ and an appropriately chosen filtration would likely shed much light on the motivic Adams spectral sequence over $\C$. Such a category might be called a ``bisynthetic''.
\end{enumerate}


%% In a different direction, the deformation of spectra associated to $E = \F_p$ and the usual notion of connectivity has also proven useful in several circumstances.
%When $p$ is odd, the study of the $\underline{\mathbb{F}}_p$-Adams spectral sequence has been held back by the fact that the 
%On the other hand, one might seek similar deformations of 
%
%\begin{rmk}
%One of our motivations for this research was the hope that a deformation-theoretic description of $\R$-motivic stable homotopy theory would suggest other profitable deformations of classical stable homotopy theories.

%It seems clear from \Cref{thm:filtered-model} that there are natural deformations to study associated to slice towers of \emph{norms} of $\MUR$.  A variant deformation at odd primes would also follow from the existence of the hypothetical $\mathrm{BP}_{\mu_p}$ spectrum.
%\end{rmk}



\subsection*{Notations and conventions}\ 

%% \todo{We should write this at the end once we know what needs to go here.}

%% \todo{Implement the hyper-modern notation package.}

Throughout this paper the term category will refer to an $\infty$-category as developed by Joyal and Lurie. In some places we use the term $1$-category, which is short-hand for 1-truncated $\infty$-category. We will also assume the reader is familiar with higher algebra as developed in \cite{HA}.
%% For an introduction to motivic spectra in this context we refer the reader to [cite].

Throughout this paper, filtered and graded objects will be ubiquitous.
We adopt the convention that a filtered object in a category $\CC$ is a diagram of the form
  \[ \cdots \to C_2 \to C_1 \to C_0 \to C_{-1} \to C_{-2} \to \cdots, \]
i.e. that the maps decrease the index variable. We provide a more complete introduction to filtered objects in \Cref{app:fil}, where we set up notation for several standard constructions.

%% Since we have totally dispensed with classical foundations we are now free to reclaim the symbol $\cong$. In this paper we will use $\cong$ when two objects are equivalent along a specified equivalence, when two objects are abstractly equivalent but no preferred equivalence is given (or they are identified along the full poly-isomorphism in the sense of [trash]) we will use the symbol $\simeq$.

Let $\CC$ denote a stable, presentably symmetric monoidal category.
We will let $\CC_p$ denote the category of $p$-complete objects in $\CC$; this category acquires a symmetric monoidal structure through the completion of the symmetric monoidal structure on $\CC$.
Similarly, if $X$ is an an object of $\CC$, we let $X_p$ denote the $p$-completion of $X$.
If $\CC$ has a set of compact generators, then we let $\CC_{ip}$ denote the ind-completion of the full subcategory of $\CC$ generated under finite colimits and retracts by the $p$-completions of compact objects. In the situation where the unit is compact and all compact objects are dualizable, this definition admits the following simplification: $\CC_{ip} \simeq \Mod ( \CC ; \o_p)$. The tensor product can then be described as the relative tensor product over $\o_p$. The reason for this equivalence is that tensoring with a dualizable object commutes with limits, so for $X$ dualizable $X_p \simeq X \otimes \o_p$.
%% the the restriction to the dualizable objects $\CC_{ip} ^{dual} \to \CC_p$ is fully faithful. Equivalently, a dualizable $\o_p$-module $M$ is already $p$-complete.

Our reason for using $\CC_{ip}$ over $\CC_p$ is that, even if the unit in $\CC$ is compact, the unit in $\CC_p$ may not be compact. On the other hand, compactness of the unit in $\CC$ implies compactness of the unit in $\CC_{ip}$. This will make certain arguments more direct in the $ip$ case.
%As a more philosophical point, the $ip$ category feels more natural as a category of $p$-adic objects since for example $\Q_p$ live in $\Ab_{ip}$ but not $\Ab_p$.
The reader who strongly prefers the usual notion of $p$-completion will be relieved to know that in convenient cases (such as those in which we work throughout this paper) $(\CC_{ip})_p \simeq \CC_p$. Therefore, all the main theorems above admit $p$-complete analogs. 
%For further discussion of aspects of $p$-completion we refer the reader to \Cref{sec:2vsi2complete}, which is logically independent from the rest of the paper (except for its treatment of $\kq$).

When $\CC$ is a stable category and $X,Y \in \CC$ are objects, we will typically let $\Map_{\CC} (X,Y)$ denote the spectrum of maps from $X$ to $Y$.
One exception is when $\CC$ comes with a natural enrichment to $C_2$-spectra, such as when $\CC = \SHR$. In that case, we will regard $\Map_{\CC} (X,Y)$ as a $C_2$-spectrum.
When we want to access the underlying space of maps, we shall use the notation $\Omega^{\infty} \Map_{\CC} (X,Y)$.

%% category and let $\CC_{ip}$ denote the category of modules over the $p$-completion of the unit.
%% Then the map $\CC \to \CC_p$ factors as $\CC \to \CC_{ip} \to \CC_p$ in $\CAlg(Pr^L)$.


%% WHAT TO SAY ABOUT ``FINITE TYPE'' OBJECTS? THIS IS IMPORTANT.

%% \todo{Recall our $p$-completion convention.}
%% \begin{dfn}
%% The category of Galois cellular $\mathbb{R}$-motivic spectra is the smallest full subcategory $\SHgc \subset \SHR$, closed under tensor products and colimits, that contains the the Picard elements $\Ss^1$, $\G_m$, $\Ss^{\C}$ and their inverses.

%% Let $\Ss_2^{\wedge}$ denote the $2$-completed unit in $\SHR$, and consider the category $\Ss_2^{\wedge}$-mod of $\Ss_2^{\wedge}$-module spectra.  We similarly define $\SHgc$ to be the smallest full subcategory of $\Ss_2^{\wedge}$-mod, closed under colimits and tensor products over $\Ss^{\wedge}_{2}$, containing the Picard elements $(\Ss^1)^{\wedge}_2, (\G_m)_2^{\wedge}$, $(\Ss^{\C})_2^{\wedge}$, and their inverses.

%% Both $\SHgc$ and $\SHgc$ are symmetric monoidal stable $\infty$-categories, and by CITE the $2$-completions of these categories are equivalent.  Note however that not every object in $\SHgc$ is $2$-complete.
%% \end{dfn}

%% \todo{At some point the notations for graded and filtered objects will have to be standardized across the paper.}


%% \todo{MANY MANY CITATIONS TO BE ADDED}

%% We will use $Y : \CC \to \CC^{\Fil}$ to denote the lax symmetric monoidal functor that sends an object $X$ to the constant filtered object at $X$. 


\subsection*{Acknowledgements}\

We thank
Tom Bachmann,
Mark Behrens,
Dan Isaksen,
Piotr Pstr\k{a}gowski and
Zhouli Xu
for helpful conversations.
During the course of this work, the middle author was supported by NSF grant DMS-1803273,
and the last author by an NSF GRFP fellowship under Grant No. 1745302.
